\newcommand{\rmcom}{\mathrm{commute}}
\newcommand{\rmcons}{\mathrm{consistency}}
\newcommand{\rmcomp}{\mathrm{completeness}}
\newcommand{\distinct}[1]{\mathsf{Distinct}_{#1}}
\newcommand{\outc}{\mathsf{Outcomes}}
\newcommand{\sfO}{\mathsf{O}}
\newcommand{\wH}{\widehat{H}}
\newcommand{\wG}{\widehat{G}}

\section{Pasting}\label{sec:ld-pasting}

\begin{theorem}[Pasting]\label{thm:ld-pasting}
  Let $(\psi, A, B, L)$ be an $(\eps, \delta, \gamma)$-good symmetric strategy for the $(m+1,q,d)$ low individual degree test.
  Let $\{G^x\}_{x \in \F_q}$ denote a set of projective sub-measurements in $\polysub{m}{q}{d}$ with the following properties:
  \begin{enumerate}
  \item (Completeness): \label{item:ld-pasting-completeness}  If $G = \E_{\bx}\sum_g G_g^{\bx}$, then
  	\begin{equation*}
	\bra{\psi} G \otimes I \ket{\psi} \geq 1 - \kappa.
	\end{equation*}
  \item (Consistency with~$A$): \label{item:ld-pasting-consistency} On average over $(\bu, \bx) \sim \F_q^{m+1}$,
	  \begin{equation*}
	  A^{u, x}_a \otimes I \simeq_{\zeta} I \otimes G^x_{[g(u)=a]}.
	  \end{equation*}
  \item (Strong self-consistency): \label{item:ld-pasting-self-consistency} On average over~$\bx \sim \F_q$,
  	\begin{equation*}
		G^x_g \otimes I \approx_{\zeta} I \otimes G^x_g.
	\end{equation*}
  \item (Boundedness): \label{item:ld-pasting-boundedness} There exists a positive-semidefinite matrix $Z^x$ for each $x \in \F_q$ such that
  	\begin{equation*}
		\E_{\bx} \bra{\psi} (I-G^{\bx})\otimes Z^{\bx} \ket{\psi} \leq \zeta
	\end{equation*}
	and for each $x \in \F_q$ and $g \in \polyfunc{m}{q}{d}$,
	\begin{equation*}
	Z^x \geq \left(\E_{\bu} A^{\bu, x}_{g(\bu)}\right).
	\end{equation*}
  \end{enumerate}
  Let $k \geq 400md$ be an integer. Let
 \begin{align*}
 \nu &= 100 k^2m \cdot \left(\eps^{1/32} + \delta^{1/32} + \gamma^{1/32} + \zeta^{1/32} + (d/q)^{1/32}\right),\\
 \sigma & = \kappa \cdot \left(1 + \frac{1}{100m}\right)
    + 2\nu + e^{- k/(80000m^2)}.
 \end{align*}
  Then there exists a ``pasted" measurement $H \in \polymeas{m+1}{q}{d}$ which satisfies the following property.
  \begin{enumerate}
  \item (Consistency with~$A$): \label{item:ld-pasting-N-consistency} On average over $\bu \sim \F_q^{m+1}$,
	  \begin{equation*}
	  A^{u}_a \otimes I \simeq_{\sigma} I \otimes H_{[h(u)=a]}.
	  \end{equation*}
  \end{enumerate}
\end{theorem}

We note that the bound we are proving is trivial when at least one of $\eps$, $\delta$, $\gamma$, $\zeta$, or $d/q$ is $\geq 1$,
as $\nu$ is at least~$1$ in that case.
Hence, we may assume that $\eps,\delta,\gamma, \zeta, d/q \leq 1$.
This will aid us when carrying out the error calculations,
as it allows us to bound terms like $\zeta^{1/2}$ by terms like $\zeta^{1/4}$.

It will be convenient to state
certain consistency relations that follow easily from the hypotheses
of the theorem.
In this section, we will  only use the axis-parallel lines test in the $(m+1)$-st direction,
		i.e.\ in the case when the line is of the form $\ell = \{(u, x) \mid x \in \F_q\}$.
		Such a line is specified by a point $u \in \F_q^m$.
		As a result, we can use the shorthand introduced in \Cref{not:conditioned-on-last-direction},
		where instead of writing $B^\ell_f$ for such a line, we write $B^u_f$.
		We will also make use of the other shorthand from \Cref{not:conditioned-on-last-direction}
		in which for a function $f:\ell \rightarrow \F_q$,
		we will sometimes write $f(x)$ instead of $f(u,x)$.
		
		We know that $(\psi, A, B, L)$ passes the axis-parallel lines test with probability $1-\eps$,
		and so it passes this test with probability $1-(m+1) \cdot \eps$ conditioned on the random 		direction $\bi \in \{1, \ldots, m+1\}$ from the test
		being equal to $m+1$.
		This means the following: if $(\bu, \bx) \sim \F_q^{m+1}$ is drawn uniformly at random and $\bell = \{(\bu, y) \mid y \in \F_q\}$
		is the line through it in the $(m+1)$-st direction, then
		\begin{equation*}
		A^{u, x}_a \ot I \simeq_{(m+1) \eps} I \ot B^{\ell}_{[f(x)=a]} = I \ot B^u_{[f(x)=a]}.
		\end{equation*}
		Hence, \Cref{prop:simeq-to-approx} and the fact that $(m+1) \leq 2m$ imply that
		\begin{equation*}
		A^{u, x}_a \ot I \approx_{4m \eps} I \ot B^u_{[f(x)=a]}.
		\end{equation*}

Next, from the assumption that the given strategy is $(\eps,
\delta, \gamma)$-good, \Cref{prop:simeq-to-approx} implies that
\begin{equation*}
I \otimes A^{u, x}_a \approx_{2\delta} A^{u, x}_a \otimes I \approx_{4m\eps} I \otimes B^{u}_{[f(x)=a]}.
\end{equation*}
As a result, \Cref{prop:triangle-inequality-for-approx_delta} implies that
\begin{equation}\label{eq:ld-abcon}
I \otimes A^{u, x}_a \approx_{8m\eps + 4\delta} I\otimes B^{u}_{[f(x)=a]}.
\end{equation}
\Cref{prop:triangle-sub} applied to \Cref{item:ld-pasting-consistency} and \Cref{eq:ld-abcon} implies
\begin{equation}\label{eq:ld-gbcon}
G^x_{[g(u) = a]} \ot I \simeq_{\nu_1} I \ot B^{u}_{[f(x)=a]},
\end{equation}
where
\begin{equation}\label{eq:ld-nu1-def}
\nu_1 = \zeta + \sqrt{8 m\eps + 4\delta}.
\end{equation}
Finally, \Cref{thm:com-main} implies that
\begin{equation}\label{eq:quote-com-main}
G^x_g G^y_h \otimes I \approx_{\nu_{\rmcom}} G^y_h G^x_g \otimes I,
\end{equation}
where
\begin{equation*}
\nu_{\rmcom} = 30 m \cdot \left( \gamma^{1/4} +\zeta^{1/4} + (d/q)^{1/4}\right).
\end{equation*}



\subsection{From measurements to sub-measurements}

Rather than designing the pasted measurement guaranteed by \Cref{thm:ld-pasting},
it is convenient to first design a pasted \emph{sub}-measurement,
and then convert it to a measurement.
The next lemma shows the bounds we achieve for this sub-measurement.

\begin{lemma}\label{lem:ld-pasting-sub-measurement}
There exists a ``pasted" sub-measurement $H \in \polysub{m+1}{q}{d}$ which satisfies the following properties.
  \begin{enumerate}
  \item (Consistency with~$A$): \label{item:ld-pasting-N-consistency-sub-measurement} On average over $\bu \sim \F_q^{m+1}$,
	  \begin{equation*}
	  A^{u}_a \otimes I \simeq_{\nu} I \otimes H_{[h(u)=a]}.
	  \end{equation*}
  \item (Completeness): \label{item:ld-pasting-N-completeness-sub-measurement} If $H =  \sum_h H_h$, then
  	\begin{equation*}
	\bra{\psi} H \otimes I \ket{\psi} \geq 1 - \kappa \cdot \left(1 + \frac{1}{100m}\right)
    - \nu - e^{- k/(80000m^2)}.
	\end{equation*}
  \end{enumerate}
\end{lemma}
\begin{proof}[Proof of \Cref{thm:ld-pasting} assuming \Cref{lem:ld-pasting-sub-measurement}]
Let~$H$ be the sub-measurement in $\polysub{m+1}{q}{d}$ guaranteed by \Cref{lem:ld-pasting-sub-measurement}.
Let $h^*$ be an arbitrary polynomial in $\polyfunc{m+1}{q}{d}$.
We define the measurement~$H_{\mathrm{meas}} \in \polymeas{m+1}{q}{d}$ as follows.
\begin{equation*}
(H_{\mathrm{meas}})_h
= \left\{\begin{array}{cl}
		H_{h^*} + (I- H) & \text{if $h = h^*$,}\\
		H_h & \text{otherwise.}
		\end{array}
		\right.
\end{equation*}
This is clearly a measurement, as the sum of its POVM elements is $H + (I -H) = H$.
In addition,
\begin{align*}
&\E_{\bu} \sum_{a \neq b} \bra{\psi} A^{\bu}_a \ot (H_{\mathrm{meas}})_{[h(\bu)=b]} \ket{\psi}\\
 ={}& \E_{\bu} \sum_{a} \sum_{h:h(\bu) \neq a} \bra{\psi} A^{\bu}_a \ot (H_{\mathrm{meas}})_{h} \ket{\psi}\\
 ={}& \E_{\bu} \sum_{a} \sum_{h:h(\bu) \neq a} \bra{\psi} A^{\bu}_a \ot H_{h} \ket{\psi}
 	+\E_{\bu} \sum_{a:h^*(\bu) \neq a}\bra{\psi} A^{\bu}_a \ot (I-H) \ket{\psi}\\
 \leq{}& \E_{\bu} \sum_{a} \sum_{h:h(\bu) \neq a} \bra{\psi} A^{\bu}_a \ot H_{h} \ket{\psi}
 	+\E_{\bu} \bra{\psi} I \ot (I-H) \ket{\psi}\\
\leq{}& (\nu) +  \left(\kappa \cdot \left(1 + \frac{1}{100m}\right)
    				+ \nu + e^{- k/(80000m^2)}\right)
					\tag{by \Cref{item:ld-pasting-N-consistency-sub-measurement,item:ld-pasting-N-completeness-sub-measurement}}\\
={}& \sigma.
\end{align*}
Thus, $A^u_a \ot I \simeq_{\sigma} I \ot (H_{\mathrm{meas}})_{[h(u)=a]}$. This completes the proof.
\end{proof}

We will now spend the rest of this section proving \Cref{lem:ld-pasting-sub-measurement},
i.e.\ designing the pasted sub-measurement~$H$.

\subsection{The pasted sub-measurement}

\begin{definition}
For $k \geq 1$, we let $\distinct{k}$ be the set of tuples $(x_1, \ldots, x_k) \in \F_q^k$ such that $x_i \neq x_j$ for all $i \neq j$.
We write $(\bx_1, \ldots, \bx_k) \sim \distinct{k}$ for a uniformly random element of this set.
\end{definition}


In manipulations involving our pasted measurement, it will often be useful to switch
the expectation over $ \distinct{k}$ with an expectation over uniformly
random $k$-tuples of elements of $\F_q$. The following proposition
bounds the distance between these two distributions.

\begin{proposition}\label{prop:ld-dnoteq}
Let $\bx = (\bx_1, \ldots, \bx_k) \sim \F_q^k$ be sampled uniformly at random
and let $\by = (\by_1, \ldots, \by_k) \sim \distinct{k}$.
Then
\begin{equation*}
d_{\mathrm{TV}}(\bx, \by) \leq \frac{k^2}{q}.
\end{equation*}
\end{proposition}
\begin{proof}
For any $z = (z_1, \ldots, z_k) \in \F_q^k$, 
\begin{align*}
\Pr[\bx = z] 
& = \frac{1}{q^k},\\
\Pr[\by = z] 
& = \left\{\begin{array}{cl}
		\frac{1}{\binom{q}{k} k!} & \text{if $z \in \distinct{k}$,}\\
		0 & \text{otherwise.}
		\end{array}\right.
\end{align*}
Because $\frac{1}{\binom{q}{k} k!} \geq \frac{1}{q^k}$, $\Pr[\bx = z] \geq \Pr[\by = z]$ if and only if $z \notin \distinct{k}$.
Hence,
\begin{align*}
d_{\mathrm{TV}}(\bx, \by)
&= \max_{S \subseteq \F_q^k}\{\Pr[\bx \in S] - \Pr[\by \in S]\}\\
&= \Pr[\bx \in \overline{\distinct{k}}] - \Pr[\by \in \overline{\distinct{k}}]
= \Pr[\bx \in \overline{\distinct{k}}].
\end{align*}
We can upper-bound this probability as follows.
\begin{align*}
\Pr[\bx \in \overline{\distinct{k}}]
&= \Pr[\exists i : \bx_i \in \{\bx_1, \ldots, \bx_{i-1}\}]\\
&\leq \sum_{i=2}^k \Pr[\bx_i \in \{\bx_1, \ldots, \bx_{i-1}\}]
\leq \sum_{i=2}^k \left(\frac{i-1}{q}\right)
= \frac{k(k-1)}{2q}.
\end{align*}
This concludes the proof.
\end{proof}

To begin, we will describe our construction of the global pasted measurement.
In fact, we will consider \emph{two} separate constructions of the global pasted measurement.
The first is given in \Cref{sec:construction-numba-one} below.
It is the more natural of the two,
but unfortunately we do not know how to prove that it works correctly.
This motivates our second construction,
given in \Cref{sec:construction-numba-two} below,
which is designed to circumvent the problems in the first construction.

\subsubsection{The first construction}\label{sec:construction-numba-one}

The first construction of $H = \{H_h\}$ is conceptually simple:
we perform the $G$ sub-measurement $d+1$ times to produce $d+1$ polynomials $g_1, \ldots, g_{d+1} \in \polyfunc{m}{q}{d}$.
We then perform polynomial interpolation to produce a single global polynomial $h \in \polyfunc{m+1}{q}{d}$.
In more detail, this construction involves three steps.

\begin{enumerate}
\item (Pasting):
Let $x_1, \ldots, x_{d+1} \in \F_q$.
We will define an initial ``sandwiched" measurement as follows:
\begin{equation*}
\widehat{H}^{x_1, \ldots, x_{d+1}}_{g_1, \ldots, g_{d+1}} = G^{x_1}_{g_1} \cdot G^{x_2}_{g_2} \cdots G^{x_{d+1}}_{g_{d+1}} \cdots G^{x_2}_{g_2} \cdot G^{x_1}_{g_1}.
\end{equation*}
$\widehat{H}^{x_1, \ldots, x_{d+1}}$ has a natural interpretation as the sub-measurement in which one performs the sub-measurements $G^{x_1}, G^{x_2}, \ldots, G^{x_{d+1}}$ one after another and outputs their results.
\item (Interpolation)
Next, let $(x_1, \ldots, x_{d+1}) \in \distinct{d+1}$.
We define the interpolated measurement
\begin{equation*}
H^{x_1, \ldots, x_{d+1}}_{h}
= \widehat{H}^{x_1, \ldots, x_{d+1}}_{h|_{x_1}, \ldots, h|_{x_{d+1}}}.
\end{equation*}
This performs the $\widehat{H}^{x_1, \ldots, x_{d+1}}$ measurement and outputs the~$h$
which is consistent with the outcomes~$g_1, \ldots, g_{d+1}$ if one exists.
(We note that this is not necessarily a sub-measurement if $(x_1, \ldots, x_{d+1}) \notin \distinct{d+1}$. This is because there may exist $h \neq h'$ for which $h|_{x_i} = (h')|_{x_i}$ for all $1 \leq i \leq d+1$.)
\item (Averaging):
Finally, we randomize over the choice of $(x_1, \ldots, x_{d+1})$. In other words, we define
\begin{equation*}
H_{h} = \E_{(\bx_1, \ldots, \bx_{d+1}) \sim \distinct{d+1}} H^{\bx_1, \ldots, \bx_{d+1}}_h.
\end{equation*}
\end{enumerate}

To show that this construction works,
we need to show two things: (i) that $H_h$ has good agreement with $A$,
and (ii) that $H$'s completeness is close to~$G$'s.
Showing (i) is simple and follows from
the approximate commutativity of $G^{x}$ and $G^{y}$ established in \Cref{thm:com-main}.
What we do not know how to do is to show (ii),
at least for general~$d$.
In fact, at first glance it even looks like it should be false!
To see why,
we note that it is possible to show that the completeness of~$H$ can be approximated as
\begin{equation*}
\bra{\psi} H \otimes I \ket{\psi} \approx \bra{\psi} G^{d+1} \otimes I \ket{\psi},
\end{equation*}
because~$H$ is performing the~$G$ measurement $d+1$ times.
We would therefore like to show that
\begin{equation}\label{eq:subtract-a-G}
\bra{\psi} G^{d+1}\otimes I \ket{\psi}
\approx \bra{\psi} G\otimes I \ket{\psi}.
\end{equation}
But if $\bra{\psi} G \otimes I \ket{\psi} = 1-\kappa$,
a ``naive analysis" would lead to the conclusion that
\begin{equation}\label{eq:dumbo-bound-for-idiots}
\bra{\psi} G^{d+1}\otimes I \ket{\psi} \approx 1 - (d+1) \cdot \kappa,
\end{equation}
which would be too great of a loss in completeness for our proof strategy to work.

However, we can show \Cref{eq:subtract-a-G} is actually correct
and therefore this ``naive analysis" is incorrect,
at least in the case of constant~$d$.
To see how this is possible,
note that \Cref{eq:dumbo-bound-for-idiots} is in fact correct when all of~$G$'s eigenvalues are equal to $1-\kappa$,
in which case $G^{d+1} = (1-\kappa)^d \cdot G$.
So we would like to show that, on the contrary,
the fact that $G$'s incompleteness is equal to $1-\kappa$
is because a ``$(1-\kappa)$ fraction" of its eigenvalues are equal to~$1$,
and the remaining ``$\kappa$ fraction" of its eigenvalues are equal to~$0$.
In this case, $G^{d+1} = G$, and so \Cref{eq:subtract-a-G} holds.

We now sketch the argument that shows this holds, at least for constant~$d$.
To do this, it is first possible to show that
\begin{equation}\label{eq:step-one-of-grand-plan}
\bra{\psi} G^{d+2} \otimes I \ket{\psi}
\approx_{\Delta} \bra{\psi} G^{d+1} \otimes I \ket{\psi},
\end{equation}
where $\Delta = \poly(m) \cdot \poly(\eps, \delta, \gamma, \zeta, d/q)$ is small.
Intuitively, this is because after performing $(d+1)$ $G$ measurements,
the outcome of another $G$ measurement is essentially determined due to interpolation.
(The proof is of this fact is slightly subtle and involves the $Z$-boundedness condition.)

From here, it is possible to derive \Cref{eq:subtract-a-G} as follows.
Write the eigendecomposition $G = \sum_i \lambda_i \ket{v_i}\bra{v_i}$ of~$G$.
This defines a probability distribution $\mu$ over eigenvectors, where eigenvector~$i$ occurs with probability $\mu(i) = \bra{\psi} (\ket{v_i}\bra{v_i} \ot I) \ket{\psi}$.
In this language, for any power~$k$ we can write
\begin{equation*}
\bra{\psi} G^k \ot I \ket{\psi}
= \bra{\psi} \Big(\sum_i \lambda_i^k \ket{v_i}\bra{v_i}\Big) \ot I \ket{\psi}
= \sum_i \lambda_i^k \cdot \mu(i)
= \E_{\bi \sim \mu}[\lambda_{\bi}^k].
\end{equation*}
Thus, \Cref{eq:step-one-of-grand-plan} is equivalent to the statement that
\begin{equation}\label{eq:equivalent-way-of-writing-grand-plan}
\Delta
\geq \bra{\psi} G^{d+1} \otimes I \ket{\psi} - \bra{\psi} G^{d+2} \otimes I \ket{\psi}
= \E_{\bi \sim \mu}[\lambda_{\bi}^{d+1}] - \E_{\bi \sim \mu}[\lambda_{\bi}^{d+2}]
= \E_{\bi \sim\mu}[\lambda_{\bi}^{d+1}(1-\lambda_{\bi})].
\end{equation}
By \Cref{eq:subtract-a-G}, our goal is to bound
\begin{equation*}
\bra{\psi} G \otimes I \ket{\psi} - \bra{\psi} G^{d+1} \otimes I \ket{\psi}
= \E_{\bi \sim \mu}[\lambda_{\bi}] - \E_{\bi \sim \mu}[\lambda_{\bi}^{d+1}]
= \E_{\bi \sim\mu}[\lambda_{\bi}(1-\lambda_{\bi}^d)].
\end{equation*}
To do so, we will use the following lemma.

\begin{lemma}\label{lem:looks-easy-but-took-me-a-while}
For any real number $0 \leq \lambda \leq 1$,
\begin{equation*}
\lambda (1-\lambda^d) \leq 2 \cdot \Big(\lambda^{d+1}(1-\lambda)\Big)^{1/(d+1)}.
\end{equation*}
\end{lemma}
\begin{proof}
The lemma is trivial for $\lambda = 1$, and so we will assume that $\lambda \neq 1$.
We note that
\begin{align*}
\lambda^{d+1} (1-\lambda^d)^{d+1}
\leq \lambda^{d+1} (1-\lambda^d)
&= \lambda^{d+1}(1-\lambda) \cdot \left(\frac{1-\lambda^d}{1-\lambda}\right)\\
&= \lambda^{d+1}(1-\lambda) \cdot (1 + \lambda + \cdots + \lambda^{d-1})
\leq d \cdot \lambda^{d+1}(1-\lambda).
\end{align*}
The lemma now follows by taking the $(d+1)$-st root of both sides and noting that $d^{1/(d+1)} \leq 2$ for all integers $d \geq 1$.
\end{proof}

Then \Cref{lem:looks-easy-but-took-me-a-while} implies that
\begin{align*}
\E_{\bi \sim\mu}[\lambda_{\bi}(1-\lambda_{\bi}^d)]
&\leq 2\cdot \E_{\bi \sim\mu}\Big[\Big(\lambda_{\bi}^{d+1}(1-\lambda_{\bi})\Big)^{1/(d+1)}\Big]\\
&\leq2\cdot \Big(\E_{\bi \sim\mu}[\lambda_{\bi}^{d+1}(1-\lambda_{\bi})]\Big)^{1/(d+1)}
							\tag{because $a \mapsto a^{1/(d+1)}$ is concave}\\
&\leq 2 \cdot \Delta^{1/(d+1)}. \tag{by \Cref{eq:equivalent-way-of-writing-grand-plan}}
\end{align*}
Thus, we have established the bound
\begin{equation*}
\bra{\psi}G \otimes I \ket{\psi}
- \bra{\psi}G^{d+1} \otimes I \ket{\psi}
\leq 2\cdot \Delta^{1/{d+1}}
= 2\cdot ( \poly(m) \cdot \poly(\eps, \delta, \gamma, \zeta,d/q))^{1/{d+1}}.
\end{equation*}
For constant~$d$, this bound is suitable for our proof, as the right-hand side is still a polynomial in the relevant parameters.
However, when~$d$ is larger, for this bound to be meaningful, we need $\eps$, $\delta$, etc. to be exponentially small in~$d$, which is a more stringent condition than we generally allow (unless, again, $d$ is a constant). 
That said, we believe this large error may be an artifact of the proof strategy
rather than something intrinsic to the construction.
We leave this to future work.



\subsubsection{The second construction}\label{sec:construction-numba-two}

The second construction of $H = \{H_h\}$ 
is designed to circumvent the problem of the first construction,
which is that its completeness was difficult to analyze.
Instead, we will design a pasted measurement
in which the ``naive analysis" actually gets us the bound we want,
which is that~$H$'s completeness is close to~$G$'s completeness.
Before describing the construction, we need the following definitions.

\begin{definition}[$G$'s incomplete part]
For each $x \in \F_q$, we write $G^x = \sum_g G^x_g$ and $G^x_{\bot} = I - G^x$
for the ``complete" and ``incomplete" parts of $G^x$, respectively.

It will be convenient to sometimes regard $G^x$ as a complete measurement
by throwing in the additional measurement outcome ``$\bot$".
To distinguish this from $G^x$ as a sub-measurement,
we will use the notation ``$\widehat{G}^x$".
In other words, we let $\widehat{G} = \{\widehat{G}^x_g\}$
be the projective measurement defined as
\begin{equation*}
\widehat{G}^x_g
= \left\{\begin{array}{rl}
	G^x_g & \text{if } g \in \polyfunc{m}{q}{d},\\
	G^x_{\bot} & \text{if } g = \bot.
	\end{array}\right.
\end{equation*}
This measurement has outcomes ranging over the set 
$\calP^+(m,q,d) := \polyfunc{m}{q}{d} \cup \{\bot\}$.
\end{definition}

\begin{definition}[Types]
A type $\tau$ is an element of $\{0, 1\}^k$ for some integer~$k$.
We write $|\tau| = \tau_1 + \cdots + \tau_k$ for the Hamming weight of~$\tau$.
We will also associate $\tau$ with the set $\{i \mid \tau_i = 1\}$ and write $i \in \tau$ if $\tau_i = 1$.
\end{definition}

Suppose we perform the $\wG$ measurement $k$ times in succession,
generating the random outcomes $\bg_1, \ldots, \bg_k$.
Let us write $\btau \in \{0, 1\}^k$ for the ``type" of these outcomes, where
\begin{equation*}
\btau_i =
\left\{\begin{array}{rl}
	1 & \text{if $\bg_i \in \polyfunc{m}{q}{d}$},\\
	0 & \text{if $\bg_i = \bot$.}
	\end{array}
	\right.
\end{equation*}
Assuming the $\bg_i$'s are not inconsistent,
then we can interpolate them to produce a global polynomial~$\bh$
whenever $|\btau| \geq d+1$.
Hence, we would like to understand the probability that $|\btau| \geq d+1$
and ensure that it is as large as possible.
The probability that the measurement $\wG$ returns a polynomial $g \in \polyfunc{m}{q}{d}$
is equal to the completeness of~$G$, which is $1-\kappa$.
This tells us that the probability that $\btau_1 = 1$ is $1-\kappa$.
We might naively expect that the same holds for the other $\btau_i$'s as well.
We might also naively expect that the $\btau_i$'s are independent.
These two assumptions should not be expected to hold in general,
as they ignore correlations between the measurements
and the fact that each measurement perturbs the state $\ket{\psi}$ for subsequent measurements to use.
However, if we make these assumptions, then we at least have a simple toy model
for the measurement outcomes: $\btau \sim \mathrm{Binomial}(k, 1-\kappa)$.

In this toy model,
we expect $|\btau| \approx k \cdot (1-\kappa)$ on average.
This was the problem with the ``naive analysis" from the first construction:
if $k = d+1$ and $\kappa$ is reasonably large (say, on the order of $1/d$),
then we don't expect $|\btau|$ to be $\geq d+1$ with high probability,
and so we can't interpolate to produce a global polynomial.
This suggests an alternative strategy:
simply choose $k$ large enough so that $k \cdot (1 - \kappa) \gg d+1$.
In fact, as we are aiming for $H$ to have completeness close to $1-\kappa$,
we should choose~$k$ so large that $|\btau| \geq d+1$ with probability 
roughly $1-\kappa$.
This is easily done with a Chernoff bound,
which is responsible for the exponential error term in \Cref{item:ld-pasting-N-completeness-sub-measurement}
of \Cref{lem:ld-pasting-sub-measurement}.
On the other hand, if we set $k$ \emph{too} large,
then we increase the risk that our $k$ outcomes $g_1, \ldots, g_k$ are inconsistent with each other,
which is an additional source of error.
This is responsible for the tradeoff between ``large" and ``small" $k$ discussed in \Cref{sec:self-improvement-and-pasting} above.


Although, this ``naive analysis" only holds in this toy model,
it still motivates our second construction of~$H$,
which we state below.
We will show that the naive analysis,
in which we treat $\btau$ as a binomial random variable
and bound $|\btau|$ using a Chernoff bound,
can actually be made formal.

\begin{definition}[The pasted measurement]
Let $k \geq d+1$ be an integer.
\begin{enumerate}
\item (Pasting):
Let $x_1, \ldots, x_k \in \F_q$.
We will define an initial ``sandwiched" measurement as follows:
\begin{equation*}
\widehat{H}^{x_1, \ldots, x_{k}}_{g_1, \ldots, g_{k}} = \widehat{G}^{x_1}_{g_1} \cdot \widehat{G}^{x_2}_{g_2} \cdots \widehat{G}^{x_{k}}_{g_{k}} \cdots \widehat{G}^{x_2}_{g_2} \cdot \widehat{G}^{x_1}_{g_1}.
\end{equation*}
\item (Interpolation):
Next, let $(x_1, \ldots, x_{k}) \in \distinct{k}$.
For any string $w \in \{0, 1\}^k$ and polynomial $h \in \polyfunc{m+1}{q}{d}$,
we define $h_w$  to be the tuple $(g_1, \ldots, g_k)  \in \calP^+(m, q, d)^k$
where $g_i = \bot$ if $w_i = 0$ and $g_i = h|_{x_i}$ otherwise.
We define the interpolated measurement
\begin{equation*}
H^{x_1, \ldots, x_{k}}_{h}
= \sum_{w : |w| \geq d+1} \widehat{H}^{x_1, \ldots, x_{k}}_{h_w}.
\end{equation*}
\item (Averaging):
Finally, we randomize over the choice of $(x_1, \ldots, x_{k})$. In other words, we define
\begin{equation*}
H_{h} = \E_{(\bx_1, \ldots, \bx_{k}) \sim \distinct{k}} H^{\bx_1, \ldots, \bx_{k}}_h.
\end{equation*}
\end{enumerate}
\end{definition}

To analyze the second construction,
we first need to show that the $\wG$ measurement
satisfies some basic properties, like commutation with itself.
These are shown in \Cref{sec:hat-consistency},
where they follow from the fact that similar properties hold for $G$.
Using this, we prove that $\widehat{H}$ is consistent with~$B$ in \Cref{sec:ld-sandwiching},
which we use to prove that $H$ is consistent with~$A$ in \Cref{sec:consistency-of-h-with-a}.
Finally, we analyze the completeness of~$H$ in \Cref{sec:completeness-of-h-low-degree}.


\subsection{Strong self-consistency and commutation of $\widehat{G}$}
\label{sec:hat-consistency}

In this section, we show that $\wG$ is strongly self-consistent and commutes with itself.
As we already know this holds for the sub-measurement~$G$,
our task essentially reduces to showing that these properties also hold for $G$'s incomplete part,
i.e.\ $G_{\perp}$.
As it is more convenient to work with $G = I - G_{\perp}$ rather than $G_{\perp}$,
we will first show that these properties hold for~$G$;
the fact that they also hold for $G_{\perp}$ will then follow as an immediate corollary.

\subsubsection{Strong self-consistency of $G_{\perp}$}

\begin{lemma}[Strong self-consistency of $G$'s complete part]\label{lem:g-complete-self-consistency}
\begin{equation*}
G^{x} \ot I \approx_{\zeta}  I \ot G^{x}.
\end{equation*}
\end{lemma}
\begin{proof}
Because~$G$ is a projective measurement, \Cref{prop:two-notions-of-self-consistency}
implies that the strong self-consistency of~$G$ from \Cref{item:ld-pasting-self-consistency} is equivalent to 
\begin{equation}\label{eq:ld-g-self-consistency}
\E_{\bx} \sum_{g} \bra{\psi} G^{\bx}_{g} \ot G^{\bx}_{g}
    \ket{\psi} \geq \E_{\bx} \sum_{g} \bra{\psi} G^{\bx}_{g} \ot
                 I\ket{\psi} - \frac{1}{2}\cdot \zeta.
\end{equation}
Our goal is to bound
\begin{align*}
&\E_{\bx} \Vert (G^{\bx}\ot I - I \ot G^{\bx}) \ket{\psi} \Vert^2\\
=~&2\cdot \E_{\bx} \bra{\psi} G^{\bx} \otimes I \ket{\psi} - 2\cdot \E_{\bx} \bra{\psi} G^{\bx} \otimes G^{\bx} \ket{\psi}\\
=~&2\cdot \E_{\bx} \sum_{g_1} \bra{\psi} G^{\bx}_{g} \otimes I \ket{\psi} - 2\cdot \E_{\bx} \sum_{g, h} \bra{\psi} G^{\bx}_{g} \otimes G^{\bx}_{h} \ket{\psi}\\
\leq~&2\cdot \E_{\bx} \sum_{g} \bra{\psi} G^{\bx}_{g} \otimes I \ket{\psi} - 2\cdot \E_{\bx} \sum_{g} \bra{\psi} G^{\bx}_{g} \otimes G^{\bx}_{g} \ket{\psi}.
\end{align*}
But this is at most~$\zeta$ by \Cref{eq:ld-g-self-consistency}.
\end{proof}

\begin{corollary}[Strong self-consistency of~$G$'s incomplete part]\label{cor:g-bot-self-consistency}
\begin{equation*}
G^{x}_{\bot} \ot I \approx_{\zeta}  I \ot G^{x}_{\bot}.
\end{equation*}
\end{corollary}
\begin{proof}
 For any~$x$,
 \begin{equation*}
 G^x_{\bot} \otimes I - I \otimes G^x_{\bot}
 = (I - G^x) \otimes I - I \otimes (I - G^x)
 = I \otimes G^x - G^x \otimes I.
 \end{equation*}
Thus,
\begin{equation*}
\E_{\bx}  \| (G^{\bx}_{\bot} \ot I -I \ot
      G^{\bx}_{\bot}) \ket{\psi} \|^2
= \E_{\bx}  \| (I \otimes G^{\bx} - G^{\bx} \otimes I) \ket{\psi} \|^2,
\end{equation*}
which is at most $\zeta$ by \Cref{lem:g-complete-self-consistency}.
\end{proof}

\subsubsection{Commutativity of $G_{\perp}$}

\begin{lemma}[Commutativity with~$G_g^x$ implies commutativity with~$G^x$]\label{lem:commutativity-switcheroo}
Let $M = \{M^x_o\}$ be a projective sub-measurement with outcomes in some set~$\calO$.
Suppose that
\begin{equation}\label{eq:M-self-consistent}
M^x_o \otimes I \approx_{\omega} I \otimes M^x_o,
\end{equation}
and
\begin{equation}\label{eq:M-commutes-with-G}
G^x_g M^y_o \otimes I \approx_{\chi} M^y_o G^x_g \otimes I
\end{equation}
over independent and uniformly random $\bx, \by \sim \F_q$. Then
\begin{equation*}
G^x M^y_o \otimes I \approx_{6\sqrt{\zeta} + 6\sqrt{\omega} + 4\sqrt{\chi}} M^y_o G^x \otimes I.
\end{equation*}
\end{lemma}
\begin{proof}
The error we wish to bound is
\begin{align}
&\E_{\bx} \sum_o \Vert (G^{\bx} M^{\by}_o - M^{\by}_o G^{\bx}) \otimes I \ket{\psi}\Vert^2\nonumber\\
=~&\E_{\bx, \by} \sum_o \bra{\psi} M^{\by}_o (G^{\bx})^2 M^{\by}_o \otimes I \ket{\psi} + \E_{\bx, \by} \sum_o \bra{\psi} G^{\bx} (M^{\by}_o)^2 G^{\bx} \otimes I \ket{\psi}\nonumber\\
& \qquad - \E_{\bx, \by} \sum_o  \bra{\psi} G^{\bx} M^{\by}_o
        G^{\bx} M^{\by}_o \ot I \ket{\psi}
        	- \E_{\bx, \by} \sum_g \bra{\psi}
        M^{\by}_o G^{\bx} M^{\by}_o G^{\bx} \ot I \ket{\psi}.\label{eq:g-commute-with-gg-error}
\end{align}
We will show that all four terms in \Cref{eq:g-commute-with-gg-error} are close to $\bra{\psi} G \otimes M \ket{\psi}$,
where $M = \E_{\by} \sum_o M^{\by}_o$.

For the first term in \Cref{eq:g-commute-with-gg-error}, we have
\begin{align*}
\E_{\bx, \by} \sum_o \bra{\psi} M^{\by}_o (G^{\bx})^2 M^{\by}_o\ot I \ket{\psi}
&= \E_{\bx, \by} \sum_o \bra{\psi} M^{\by}_o G^{\bx} M^{\by}_o\ot I \ket{\psi} \tag{because~$G$ is projective}\\
&= \E_{\by}\sum_o \bra{\psi} M^{\by}_o G M^{\by}_o \ot I \ket{\psi}\\
&\approx_{2\sqrt{\omega}} \E_{\by}\sum_o \bra{\psi} G \ot M^{\by}_o  \ket{\psi} \tag{by \Cref{prop:switch-sandwich} and \Cref{eq:M-self-consistent}}\\
&= \bra{\psi} G \ot M \ket{\psi}.
\end{align*}
Similarly, for the second term in \Cref{eq:g-commute-with-gg-error}, we have
\begin{align*}
\E_{\bx, \by} \sum_o \bra{\psi} G^{\bx} (M^{\by}_o)^2 G^{\bx} \otimes I \ket{\psi}
&=\E_{\bx, \by} \sum_o \bra{\psi} G^{\bx} M^{\by}_o G^{\bx} \otimes I \ket{\psi} \tag{because~$M$ is projective}\\
&=\E_{\bx}  \bra{\psi} G^{\bx} M G^{\bx} \otimes I \ket{\psi}\\
&\approx_{2\sqrt{\zeta}} \E_{\bx} \bra{\psi} M \ot G^{\bx}  \ket{\psi} \tag{by \Cref{prop:switch-sandwich} and \Cref{lem:g-complete-self-consistency}}\\
&= \bra{\psi} M \ot G \ket{\psi}.
\end{align*}

For the third term in \Cref{eq:g-commute-with-gg-error}, we begin by claiming that
\begin{align}
\E_{\bx, \by} \sum_o \bra{\psi} G^{\bx} M^{\by}_o
        G^{\bx} M^{\by}_o \ot I \ket{\psi}
& = 
\E_{\bx, \by} \sum_{o, g}  \bra{\psi} G^{\bx} M^{\by}_o
        G^{\bx}_{g} M^{\by}_o \ot I \ket{\psi}\nonumber\\
& = 
\E_{\bx, \by} \sum_{o, g}  \bra{\psi} G^{\bx} M^{\by}_o G^{\bx}_g
        G^{\bx}_{g} M^{\by}_o \ot I \ket{\psi} \tag{because~$G$ is projective}\\
& \approx_{\sqrt{\chi}} 
\E_{\bx, \by} \sum_{o, g}  \bra{\psi}  G^{\bx} M^{\by}_o G^{\bx}_g
        M^{\by}_o G^{\bx}_{g}  \ot I \ket{\psi}.\label{eq:split-G-and-commute}
\end{align}
To show this, we bound the magnitude of the difference.
\begin{multline*}
\Big|\E_{\bx, \by} \sum_{o, g}  \bra{\psi}  (G^{\bx}
        M^{\by}_{o} G^{\bx}_g \ot I) \cdot ((G^{\bx}_g M^{\by}_o - M^{\by}_o G^{\bx}_{g}) \otimes I)\ket{\psi}\Big|\\
\leq \sqrt{\E_{\bx, \by} \sum_{o, g}  \bra{\psi} (G^{\bx} M^{\by}_o G^{\bx}_g
        M^{\by}_{o} G^{\bx}) \ot I\ket{\psi}}\\
\cdot \sqrt{\E_{\bx, \by} \sum_{o, g}  \bra{\psi}  ((M^{\by}_o G^{\bx}_g  -  G^{\bx}_{g} M^{\by}_o) \cdot (G^{\bx}_g M^{\by}_o - M^{\by}_o G^{\bx}_{g}) \otimes I) \ket{\psi}}.
\end{multline*}
The expression inside the first square root is at most~$1$ because~$G$ and~$M$ are sub-measurements.
The expression inside the second square root is at most $\chi$ by \Cref{eq:M-commutes-with-G}.
Next, we claim that
\begin{equation}\label{eq:move-G-for-great-justice}
\eqref{eq:split-G-and-commute}
\approx_{\sqrt{\zeta}} \E_{\bx, \by} \sum_{o, g}  \bra{\psi}   G^{\bx} M^{\by}_o
        G^{\bx}_{g} M^{\by}_o \ot G^{\bx}_g\ket{\psi}.
\end{equation}
To show this, we bound the magnitude of the difference.
\begin{align*}
&\Big|  \E_{\bx, \by} \sum_{o, g}  \bra{\psi} (G^{\bx} M^{\by}_{o} G^{\bx}_g \ot I)
	\cdot ((M^{\by}_{o} \otimes I) \cdot(G^{\bx}_g \otimes I - I \otimes G^{\bx}_g))\ket{\psi}\Big|\\
&\leq  \sqrt{\E_{\bx, \by} \sum_{o, g} \bra{\psi} (G^{\bx} M^{\by}_o G^{\bx}_g M^{\by}_{o} G^{\bx}) \ot I \ket{\psi}}\\
&\quad \cdot \sqrt{ \E_{\bx, \by} \sum_{o, g}  \bra{\psi}   ((G^{\bx}_g \otimes I - I \otimes G^{\bx}_g) \cdot (M^{\by}_{o} \otimes I) \cdot (G^{\bx}_g \otimes I - I \otimes G^{\bx}_g)) \ket{\psi}}.
\end{align*}
The expression inside the first square root is at most~$1$ because~$G$ and~$M$ are sub-measurements.
The expression inside the second square root is
\begin{align*}
& \E_{\bx}\sum_g \bra{\psi}(G^{\bx}_g \ot I - I \ot G^{\bx}_g) \cdot\Big(\E_{\by} \sum_o M^{\by}_{o} \ot I\Big) \cdot (G^{\bx}_g \ot I - I \ot G^{\bx}_g) \ket{\psi}\\
\leq~&\E_{\bx}\sum_g \bra{\psi}(G^{\bx}_g \ot I - I \ot G^{\bx}_g)^2 \ket{\psi}. \tag{because~$M$ is a sub-measurement}
\end{align*} 
This is at most~$\zeta$ by \Cref{item:ld-pasting-self-consistency}.
Next, we claim that
\begin{equation}\label{eq:move-G-for-even-greater-justice}
\eqref{eq:move-G-for-great-justice}
\approx_{\sqrt{\zeta}} \E_{\bx, \by} \sum_{o, g}  \bra{\psi}  G^{\bx}_g  G^{\bx} M^{\by}_o
        G^{\bx}_{g} M^{\by}_o \ot I\ket{\psi}.
\end{equation}
To show this, we bound the magnitude of the difference.
\begin{align*}
&\Big|  \E_{\bx, \by} \sum_{o, g}  \bra{\psi} ((G^{\bx}_g \ot I - I \ot G^{\bx}_g) \cdot (G^{\bx} M^{\by}_o \otimes I)) 
	\cdot  (G^{\bx}_{g} M^{\by}_o \otimes I) 
        \ket{\psi}\Big|\\
&\leq \sqrt{ \E_{\bx, \by} \sum_{o, g}  \bra{\psi} ((G^{\bx}_g \ot I - I \ot G^{\bx}_g) \cdot (G^{\bx} M^{\by}_o G^{\bx} \ot I) \cdot (G^{\bx}_g \ot I - I \ot G^{\bx}_g))\ket{\psi}}\\
&\quad \cdot \sqrt{ \E_{\bx, \by} \sum_{o, g}  \bra{\psi}  (M^{\by}_{o} G^{\bx}_g M^{\by}_o) \otimes I \ket{\psi}}.
\end{align*}
The expression inside the first square root is
\begin{align*}
& \E_{\bx} \sum_{g}  \bra{\psi} ((G^{\bx}_g \ot I - I \ot G^{\bx}_g) \cdot \Big(\E_{\by} \sum_o G^{\bx} M^{\by}_o G^{\bx} \ot I\Big) \cdot (G^{\bx}_g \ot I - I \ot G^{\bx}_g))\ket{\psi}\\
\leq~&\E_{\bx} \sum_{g}  \bra{\psi} (G^{\bx}_g \ot I - I \ot G^{\bx}_g)^2\ket{\psi}. \tag{because~$G$ and~$M$ are sub-measurements}
\end{align*} 
This is at most~$\zeta$ by \Cref{item:ld-pasting-self-consistency}.
The expression inside the second square root is at most~$1$ because~$G$ and~$M$ are sub-measurements.
Next, we claim that
\begin{align}
\eqref{eq:move-G-for-even-greater-justice}
& = \E_{\bx, \by} \sum_{o, g}  \bra{\psi}   G^{\bx}_{g} M^{\by}_o
        G^{\bx}_{g} M^{\by}_o \ot I\ket{\psi} \tag{because~$G$ is projective}\\
& \approx_{\sqrt{\chi}} 
\E_{\bx, \by} \sum_{o,g}  \bra{\psi}   M^{\by}_{o} G^{\bx}_g G^{\bx}_g
        M^{\by}_{o}  \ot I\ket{\psi}. \label{eq:commute-the-G-yet-again}
\end{align}
To show this, we bound the magnitude of the difference.
\begin{multline*}
\Big|\E_{\bx, \by} \sum_{o, g}  \bra{\psi} ((G^{\bx}_g
        M^{\by}_{o} - M^{\by}_o G^{\bx}_g)  \ot I)  \cdot(G^{\bx}_{g} M^{\by}_o  \ot I ) \ket{\psi}\Big|\\
\leq\sqrt{\E_{\bx, \by} \sum_{o, g} \bra{\psi} ((G^{\bx}_g M^{\by}_o - M^{\by}_o G^{\bx}_g) \cdot (M^{\by}_o
        G^{\bx}_{g} - G^{\bx}_g M^{\by}_o)  \ot I)\ket{\psi}}\\
        \cdot \sqrt{\E_{\bx, \by} \sum_{o, g}  \bra{\psi} (M^{\by}_o G^{\bx}_g M^{\by}_o) \otimes I \ket{\psi}}.
\end{multline*}
The expression inside the first square root is at most $\chi$ by \Cref{eq:M-commutes-with-G}.
The expression inside the second square root is at most~$1$ because~$G$ and~$M$ are sub-measurements.
Finally,
\begin{align*}
\eqref{eq:commute-the-G-yet-again}
&=\E_{\bx, \by} \sum_{o, g}  \bra{\psi}   M^{\by}_{o} G^{\bx}_g
        M^{\by}_{o}  \ot I\ket{\psi} \tag{because~$G$ is projective}\\
&=\E_{\by} \sum_o \bra{\psi} M^{\by}_{o} G M^{\by}_{o} \otimes I \ket{\psi}\\
&\approx_{2\sqrt{\omega}} \E_{\by} \sum_o \bra{\psi} G \ot M^{\by}_o  \ket{\psi} \tag{by \Cref{prop:switch-sandwich} and \Cref{eq:M-self-consistent}}\\
&= \bra{\psi} G \ot M \ket{\psi}.
\end{align*}
In total, this shows that
\begin{equation}\label{eq:term-three-for-use-right-now}
\E_{\bx, \by} \sum_o  \bra{\psi} G^{\bx} M^{\by}_o
        G^{\bx} M^{\by}_o \ot I \ket{\psi}
        \approx_{2 \sqrt{\zeta} + 2 \sqrt{\omega} + 2 \sqrt{\chi}} \bra{\psi} G \ot M \ket{\psi}.
\end{equation}
The fourth term in \Cref{eq:g-commute-with-gg-error} is the Hermitian conjugate of the third term. As a result, \Cref{eq:term-three-for-use-right-now} implies that
\begin{equation*}
\E_{\bx, \by} \sum_o  \bra{\psi} M^{\by}_o G^{\bx} M^{\by}_o G^{\bx} \ot I \ket{\psi}
        \approx_{2 \sqrt{\zeta} + 2 \sqrt{\omega} + 2 \sqrt{\chi}} \bra{\psi} G \ot M \ket{\psi}
\end{equation*}
as well.

In total, this gives an error of
\begin{equation*}
2\sqrt{\omega}
+ 2\sqrt{\zeta}
+ 2 \cdot\left(2 \sqrt{\zeta} + 2 \sqrt{\omega} + 2 \sqrt{\chi}\right)
= 6 \sqrt{\zeta} + 6 \sqrt{\omega} + 4 \sqrt{\chi}.
\end{equation*}
This proves the claimed bound.
\end{proof}

\begin{corollary}[Commutativity of~$G$'s complete part]\label{cor:commuting-with-G-complete}
The following commutation relations hold.
\begin{align*}
G^x_g  G^y \ot I &\approx_{\nu_2}  G^y G^x_g \otimes I\\
G^x G^y \ot I &\approx_{\nu_2}  G^y G^x \otimes I,
\end{align*}
where
\begin{equation*}
\nu_2 = 36 m \cdot \left(\gamma^{1/16} +  \zeta^{1/16} + (d/q)^{1/16}\right).
\end{equation*}
\end{corollary}
\begin{proof}
By \Cref{eq:quote-com-main},
\begin{equation}\label{eq:com-main-copy}
G^x_g G^y_h \otimes I \approx_{\nu_{\rmcom}} G^y_h G^x_g \otimes I.
\end{equation}
Now we apply \Cref{lem:commutativity-switcheroo} to \Cref{eq:com-main-copy}.
To do so, we set the ``$\{M^y_o\}$" sub-measurement to be $\{G^x_g\}$,
and therefore $\calO = \polyfunc{m}{q}{d}$.
This implies that
\begin{equation}\label{eq:applied-the-lemma-to-com-main-copy}
G^x_g G^y \otimes I \approx_{\theta_1} G^y G^x_g \otimes I
\end{equation}
for $\theta_1 = 12 \sqrt{\zeta} + 4 \sqrt{\nu_{\rmcom}}$.

Next, we apply \Cref{lem:commutativity-switcheroo} to \Cref{eq:applied-the-lemma-to-com-main-copy}.
This time, we let $\calO$ be a set containing a single outcome,
and for this outcome~$o$, we set $M^y_o = G^y$.
This implies that
\begin{equation*}
G^x G^y \otimes I \approx_{\theta_2} G^y G^x \otimes I
\end{equation*}
for $\theta_2 = 12 \sqrt{\zeta} + 4 \sqrt{\theta_1}$.
This uses \Cref{lem:g-complete-self-consistency} for the strong self-consistency of $\{M^y_o\}$.

We now show that $\nu_2$ bounds $\theta_1$ and $\theta_2$.
First, using $\sqrt{30} \leq 6$, we have
\begin{align*}
\theta_1 = 12 \sqrt{\zeta} + 4 \sqrt{\nu_{\rmcom}}
&= 12 \sqrt{\zeta} + 4 \sqrt{30 m \cdot \left(\gamma^{1/4} + \zeta^{1/4} + (d/q)^{1/4}\right)}\\
&\leq 12 \zeta^{1/8} + 24m \cdot \left(\gamma^{1/8} +\zeta^{1/8} +  (d/q)^{1/8}\right)\\
&\leq 36 m \cdot \left(\gamma^{1/8} + \zeta^{1/8} +  (d/q)^{1/8}\right).
\end{align*}
This is clearly less than $\nu_2$.
Next,  we have
\begin{align*}
\theta_2 = 12 \sqrt{\zeta} + 4 \sqrt{\theta_1}
&\leq 12 \sqrt{\zeta} + 4 \sqrt{36 m \cdot \left(\gamma^{1/8} + \zeta^{1/8} +  (d/q)^{1/8}\right)}\\
&\leq 12 \zeta^{1/16} + 24m \cdot \left(\gamma^{1/16} + \zeta^{1/16} + (d/q)^{1/16}\right)\\
&\leq 36 m \cdot \left( \gamma^{1/16} + \zeta^{1/16} + (d/q)^{1/16}\right).
\end{align*}
This is equal to $\nu_2$, which completes the proof.
\end{proof}

\begin{corollary}[Commutativity of~$G$'s incomplete part]\label{cor:commuting-with-G-incomplete}
The following commutation relations hold.
\begin{align*}
G^x_g  G^y_{\bot} \ot I &\approx_{\nu_2}  G^y_{\bot} G^x_g \otimes I\\
G^x_{\bot} G^y_{\bot} \ot I &\approx_{\nu_2}  G^y_{\bot} G^x_{\bot} \otimes I,
\end{align*}
where $\nu_2$ is as in \Cref{cor:commuting-with-G-complete}.
\end{corollary}
\begin{proof}
First, we note that
\begin{equation*}
G^x_g  G^y_{\bot} - G^y_{\bot} G^x_g
= G^x_g  \cdot (I - G^y) - (I- G^y)\cdot G^x_g
= G^y G^x_g - G^x_g G^y.
\end{equation*}
Hence, by \Cref{cor:commuting-with-G-complete},
\begin{equation*}
\E_{\bx, \by}\sum_g \Vert (G^{\bx}_g  G^{\by}_{\bot} - G^{\by}_{\bot} G^{\bx}_g) \otimes I \ket{\psi} \Vert^2
= \E_{\bx, \by}\sum_g \Vert (G^{\by} G^{\bx}_g - G^{\bx}_g G^{\by}) \otimes I \ket{\psi} \Vert^2
\leq \nu_2.
\end{equation*}
Next, we note that
\begin{align*}
G^x_{\bot} G^y_{\bot}
- G^y_{\bot} G^x_{\bot}
&= (I - G^x) \cdot (I - G^y) - (I - G^y) \cdot (I - G^x)\\
& = (I - G^x - G^y + G^x G^y) - (I - G^y - G^x + G^y G^x)\\
&= G^x G^y - G^y G^x.
\end{align*}
As a result, by \Cref{cor:commuting-with-G-complete},
\begin{equation*}
\E_{\bx, \by} \Vert (G^{\bx}_{\bot}  G^{\by}_{\bot} - G^{\by}_{\bot} G^{\bx}_{\bot}) \otimes I \ket{\psi} \Vert^2
= \E_{\bx, \by} \Vert (G^{\by} G^{\bx} - G^{\bx} G^{\by}) \otimes I \ket{\psi} \Vert^2
\leq \nu_2.\qedhere
\end{equation*}
\end{proof}

\subsubsection{Putting everything together}

Now we combine the results of the previous two sections
to show our strong self-consistency and commutation results for~$\wG$.

\begin{corollary}[Strong self-consistency and commutation of~$\wG$]\label{cor:G-hat-facts}
  $\widehat{G}$ obeys the
  following strong self-consistency and commutation properties.
  \begin{align}
    \widehat{G}^{x}_{g} \ot I &\approx_{2\zeta} I \ot \widehat{G}^{x}_{g}, \label{eq:gselfconall}\\
    \widehat{G}^{x}_{g}\widehat{G}^{y}_{h} \ot I &\approx_{\nu_3} \widehat{G}^{y}_{h}
                               \widehat{G}^{x}_{g} \ot I, \label{eq:gcomall}
  \end{align}
  where
  \[ \nu_{3} = 138 m \cdot \left(\zeta^{1/16} + \gamma^{1/16} + (d/q)^{1/16}\right). \]
\end{corollary}
\begin{proof}
We begin with \Cref{eq:gselfconall}. To prove this, we wish to bound
\begin{align*}
&\E_{\bx} \sum_{g} \Vert(\widehat{G}^{\bx}_g \otimes I - I \otimes \widehat{G}^{\bx}_g) \ket{\psi} \Vert^2\\
=~&\E_{\bx} \sum_{g \in \polyfunc{m}{q}{d}} \Vert(G^{\bx}_g \otimes I - I \otimes G^{\bx}_g) \ket{\psi} \Vert^2
	+ \E_{\bx} \Vert(G^{\bx}_{\bot} \otimes I - I \otimes G^{\bx}_{\bot}) \ket{\psi} \Vert^2\\
\leq~& \zeta + \zeta = 2\zeta,
\end{align*}
by \Cref{item:ld-pasting-self-consistency} and \Cref{cor:g-bot-self-consistency}.

Next, we show \Cref{eq:gcomall}.
To prove this, we wish to bound
\begin{align*}
&\E_{\bx, \by} \sum_{g, h} \Vert(\widehat{G}^{\bx}_g \widehat{G}^{\by}_h - \widehat{G}^{\by}_h \widehat{G}^{\bx}_g) \otimes I \ket{\psi}\Vert^2\\
=~&\E_{\bx, \by} \sum_{g, h \in \polyfunc{m}{q}{d}} \Vert(G^{\bx}_g G^{\by}_h - G^{\by}_h G^{\bx}_g) \otimes I \ket{\psi}\Vert^2
	+ \E_{\bx, \by} \Vert(G^{\bx}_{\bot} G^{\by}_{\bot} - G^{\by}_{\bot} G^{\bx}_{\bot}) \otimes I \ket{\psi}\Vert^2\\
& \quad + \E_{\bx, \by} \sum_{g\in \polyfunc{m}{q}{d}} \Vert(G^{\bx}_g G^{\by}_{\bot} - G^{\by}_{\bot} G^{\bx}_g) \otimes I \ket{\psi}\Vert^2
	+ \E_{\bx, \by} \sum_{h\in \polyfunc{m}{q}{d}} \Vert(G^{\bx}_{\bot} G^{\by}_h - G^{\by}_h G^{\bx}_{\bot}) \otimes I \ket{\psi}\Vert^2\\
\leq~& \nu_{\rmcom} + 3\nu_2,
\end{align*}
by \Cref{eq:quote-com-main} and \Cref{cor:commuting-with-G-incomplete}.
We can therefore bound this by
\begin{align*}
\nu_{\rmcom} + 3\nu_2
&= 30 m \cdot \left(\gamma^{1/4} +\zeta^{1/4} +  (d/q)^{1/4}\right)
	+ 3 \cdot 36 m \cdot \left(\gamma^{1/16} + \zeta^{1/16} +  (d/q)^{1/16}\right)\\
&\leq 30 m \cdot \left(\gamma^{1/16} + \zeta^{1/16} + (d/q)^{1/16}\right)
	+ 3 \cdot 36 m \cdot \left(\gamma^{1/16} + \zeta^{1/16} +  (d/q)^{1/16}\right)\\
&= 138 m \cdot \left( \gamma^{1/16} + \zeta^{1/16} +(d/q)^{1/16}\right).	
\end{align*}
This completes the proof.
\end{proof}


\subsection{Consistency of the sandwich $\widehat{H}$ with~$B$}
\label{sec:ld-sandwiching}

In the next lemmas, we will show that $\widehat{H}$ is consistent with the lines
measurement $B$. To start, we show some self-consistency and
commutativity properties of $\widehat{H}$.



\begin{lemma}[Commuting past multiple $\widehat{G}$'s]
\label{lem:commute-g-half-sandwich}
For all $k \geq 2$,
  \[ \widehat{G}^{x_1}_{g_1} \widehat{G}^{x_2}_{g_2} \cdots \widehat{G}^{x_k}_{g_k} \ot I \approx_{\nu_4}
    \widehat{G}^{x_2}_{g_2} \cdots \widehat{G}^{x_k}_{g_k} \widehat{G}^{x_1}_{g_1} \ot I, \]
where
\begin{equation*}
\nu_4 = 426 k^2 m \cdot \left(\gamma^{1/16} +\zeta^{1/16} +  (d/q)^{1/16}\right).
\end{equation*}
\end{lemma}
\begin{proof}
This proof will consist of multiple applications of \Cref{eq:gselfconall,eq:gcomall};
for each line, we will specify which equation to apply.
Each line will also involve an application of \Cref{prop:cab-approx-delta},
which we will specify only implicitly.
\begin{align*}
&  \widehat{G}^{x_1}_{g_1} \widehat{G}^{x_2}_{g_2} \cdots \widehat{G}^{x_k}_{g_k} \ot I\\
\approx_{2\zeta}~& \widehat{G}^{x_1}_{g_1} \widehat{G}^{x_2}_{g_2} \cdots \widehat{G}^{x_{k-1}}_{g_{k-1}} \ot \widehat{G}^{x_k}_{g_k} \tag{by \Cref{eq:gselfconall}}\\
\cdots~&\\
\approx_{2\zeta}~& \widehat{G}^{x_1}_{g_1} \widehat{G}^{x_2}_{g_2} \ot \widehat{G}^{x_k}_{g_k} \cdots \widehat{G}^{x_{3}}_{g_{3}}\tag{by \Cref{eq:gselfconall}}\\
\approx_{\nu_{3}}~&  \widehat{G}^{x_2}_{g_2}G^{x_1}_{g_1} \ot \widehat{G}^{x_k}_{g_k} \cdots \widehat{G}^{x_{3}}_{g_{3}}\tag{by \Cref{eq:gcomall}}\\
\approx_{2\zeta}~&  \widehat{G}^{x_2}_{g_2}\widehat{G}^{x_1}_{g_1} \widehat{G}^{x_3}_{g_3} \ot \widehat{G}^{x_k}_{g_k} \cdots \widehat{G}^{x_{4}}_{g_{4}}\tag{by \Cref{eq:gselfconall}}\\
\approx_{\nu_{3}}~&  \widehat{G}^{x_2}_{g_2}\widehat{G}^{x_3}_{g_3} \widehat{G}^{x_1}_{g_1} \ot \widehat{G}^{x_k}_{g_k} \cdots \widehat{G}^{x_{4}}_{g_{4}}\tag{by \Cref{eq:gcomall}}\\
\cdots~&\\
\approx_{2\zeta}~&  \widehat{G}^{x_2}_{g_2}\cdots \widehat{G}^{x_{k-1}}_{g_{k-1}} \widehat{G}^{x_1}_{g_1} \widehat{G}^{x_k}_{g_k} \ot I\tag{by \Cref{eq:gselfconall}}\\
\approx_{\nu_{3}}~&  \widehat{G}^{x_2}_{g_2}\cdots \widehat{G}^{x_{k-1}}_{g_{k-1}} \widehat{G}^{x_k}_{g_k} \widehat{G}^{x_1}_{g_1} \ot I.\tag{by \Cref{eq:gcomall}}
\end{align*}
In total, we have $(k-2) + (k-2) \leq 2k$ applications of \Cref{eq:gselfconall} with error $2\zeta$ each
and $(k-1) \leq k$ applications of \Cref{eq:gcomall} with error $\nu_3$ each.
By \Cref{prop:triangle-inequality-for-approx_delta}, this implies that 
\begin{equation*}
G^{x_1}_{g_1} G^{x_2}_{g_2} \cdots G^{x_k}_{g_k} \ot I
\approx_{3k \cdot (4k \zeta + k\nu_3)} G^{x_2}_{g_2}\cdots G^{x_{k-1}}_{g_{k-1}} G^{x_k}_{g_k} G^{x_1}_{g_1} \ot I,
\end{equation*}
as claimed.
We can bound this as follows.
\begin{align*}
3k \cdot (4k \zeta + k\nu_3)
& = 12k^2 \zeta + 3k^2 \nu_3\\
&= 12k^2 \zeta + 3 k^2 \cdot 138 m \cdot \left(\gamma^{1/16} + \zeta^{1/16} + (d/q)^{1/16}\right)\\
&\leq 12k^2 \zeta^{1/16} + 414 k^2 m \cdot \left(\gamma^{1/16} + \zeta^{1/16} + (d/q)^{1/16}\right)\\
&\leq 426 k^2 m \cdot \left(\gamma^{1/16} + \zeta^{1/16} + (d/q)^{1/16}\right).
\end{align*}
This completes the proof.
\end{proof}

\begin{lemma}[Consistency of $\widehat{H}$ with $B$]
\label{lem:ld-sandwich-line-one-point}
For any $1 \leq i \leq k$,
\begin{equation*}
\E_{\bu}  \E_{\bx_1, \dots, \bx_k} \sum_{g_1, \dots, g_k: g_i \neq \bot}  \sum_{a \neq
    g_i(\bu)} \bra{\psi}
  \widehat{H}^{\bx_1, \dots, \bx_k}_{g_1, \dots, g_k} \ot B^{\bu}_{[f(\bx_i) =
    a]} \ket{\psi}  \leq \nu_5,
 \end{equation*}
where
\begin{equation*}
    \nu_5 = 43 km \cdot \left(\eps^{1/32} + \delta^{1/32} + \gamma^{1/32} + \zeta^{1/32} + (d/q)^{1/32}\right).
\end{equation*}
\end{lemma}
\begin{proof}
To begin, we note that
\begin{align*}
\sum_{g_{i+1}, \ldots, g_k} \widehat{H}^{x_1, \ldots, x_k}_{g_1, \ldots, g_k}
& = \sum_{g_{i+1}, \ldots, g_k} \widehat{G}^{x_1}_{g_1} \cdots \widehat{G}^{x_k}_{g_k} \cdots \widehat{G}^{x_1}_{g_1}\\
& = \sum_{g_{i+1}, \ldots, g_{k-1}} \widehat{G}^{x_1}_{g_1} \cdots \Big(\sum_{g_k} \widehat{G}^{x_k}_{g_k} \Big)\cdots \widehat{G}^{x_1}_{g_1}\\
& = \sum_{g_{i+1}, \ldots, g_{k-1}} \widehat{G}^{x_1}_{g_1} \cdots \widehat{G}^{x_{k-1}}_{g_{k-1}} \cdot I \cdot \widehat{G}^{x_{k-1}}_{g_{k-1}}\cdots \widehat{G}^{x_1}_{g_1} \tag{because $\widehat{G}^{x_k}$ is a measurement}\\
& = \sum_{g_{i+1}, \ldots, g_{k-1}} \widehat{G}^{x_1}_{g_1} \cdots \widehat{G}^{x_{k-1}}_{g_{k-1}} \cdots \widehat{G}^{x_1}_{g_1} \tag{because $\widehat{G}^{x_{k-1}}$ is projective}\\
& \cdots \\
& = \widehat{G}^{x_1}_{g_1} \cdots \widehat{G}^{x_{i}}_{g_{i}} \cdots \widehat{G}^{x_1}_{g_1}\\
& = \widehat{H}^{x_1, \ldots, x_i}_{g_1, \ldots, g_i}.
\end{align*}
As a result,
\begin{align}
&\E_{\bu}  \E_{\bx_1, \dots, \bx_k} \sum_{g_1, \dots, g_k: g_i \neq \bot}  \sum_{a \neq
    g_i(\bu)} \bra{\psi}
  \widehat{H}^{\bx_1, \dots, \bx_k}_{g_1, \dots, g_k} \ot B^{\bu}_{[f(\bx_i) =
    a]} \ket{\psi}\nonumber\\
=~&\E_{\bu}  \E_{\bx_1, \dots, \bx_i} \sum_{g_1, \dots, g_i: g_i \neq \bot}  \sum_{a \neq
    g_i(\bu)} \bra{\psi}
  \widehat{H}^{\bx_1, \dots, \bx_i}_{g_1, \dots, g_i} \ot B^{\bu}_{[f(\bx_i) =
    a]} \ket{\psi}.\label{eq:delete-extraneous-coordinates}
\end{align}
For shorthand, we write
\begin{equation*}
\widehat{G}^{x_{<i}}_{g_{<i}} = \widehat{G}^{x_1}_{g_1} \cdots \widehat{G}^{x_{i-1}}_{g_{i-1}}.
\end{equation*}
Then we claim that
\begin{align}
\eqref{eq:delete-extraneous-coordinates}
& = \E_{\bu}  \E_{\bx_1, \dots, \bx_i} \sum_{g_1, \dots, g_i: g_i \neq \bot}   \bra{\psi}
  \widehat{G}^{\bx_{<i}}_{g_{<i}}  \cdot \widehat{G}^{\bx_i}_{g_i} \cdot \widehat{G}^{\bx_i}_{g_i} \cdot (\widehat{G}^{\bx_{<i}}_{g_{<i}} )^\dagger \ot (I - B^{\bu}_{[f(\bx_i) =
    g_i(\bu)]}) \ket{\psi}\nonumber\\
& \approx_{\sqrt{\nu_4}} \E_{\bu}  \E_{\bx_1, \dots, \bx_i} \sum_{g_1, \dots, g_i: g_i \neq \bot} \bra{\psi}\widehat{G}^{\bx_i}_{g_i} \cdot
  \widehat{G}^{\bx_{<i}}_{g_{<i}}  \cdot  \widehat{G}^{\bx_i}_{g_i}\cdot (\widehat{G}^{\bx_{<i}}_{g_{<i}} )^\dagger \ot (I - B^{\bu}_{[f(\bx_i) =
    g_i(\bu)]}) \ket{\psi}.\label{eq:gonna-need-a-bigger-cauchy-schwarz}
\end{align}
To show this, we bound the magnitude of the difference.
\begin{multline*}
\E_{\bu}  \E_{\bx_1, \dots, \bx_i} \sum_{g_1, \dots, g_i: g_i \neq \bot}  \bra{\psi}
  ((\widehat{G}^{\bx_{<i}}_{g_{<i}}  \cdot \widehat{G}^{\bx_i}_{g_i} - \widehat{G}^{\bx_i}_{g_i}  \cdot \widehat{G}^{\bx_{<i}}_{g_{<i}}) \otimes I) \cdot (\widehat{G}^{\bx_i}_{g_i} \cdot (\widehat{G}^{\bx_{<i}}_{g_{<i}} )^\dagger \ot (I - B^{\bu}_{[f(\bx_i) =
    g_i(\bu)]})) \ket{\psi}\\
    \leq 
    \sqrt{\E_{\bu}  \E_{\bx_1, \dots, \bx_i} \sum_{g_1, \dots, g_i: g_i \neq \bot}   \bra{\psi}
  ((\widehat{G}^{\bx_{<i}}_{g_{<i}}  \cdot \widehat{G}^{\bx_i}_{g_i} - \widehat{G}^{\bx_i}_{g_i}  \cdot \widehat{G}^{\bx_{<i}}_{g_{<i}})  \cdot  (\widehat{G}^{\bx_i}_{g_i}\cdot (\widehat{G}^{\bx_{<i}}_{g_{<i}})^\dagger - (\widehat{G}^{\bx_{<i}}_{g_{<i}})^\dagger \cdot \widehat{G}^{\bx_i}_{g_i})) \otimes I \ket{\psi}}\\
  \cdot \sqrt{\E_{\bu}  \E_{\bx_1, \dots, \bx_i} \sum_{g_1, \dots, g_i: g_i \neq \bot}  \bra{\psi}
  (\widehat{G}^{\bx_{<i}}_{g_{<i}}  \cdot \widehat{G}^{\bx_i}_{g_i} \cdot (\widehat{G}^{\bx_{<i}}_{g_{<i}} )^\dagger) \ot (I - B^{\bu}_{[f(\bx_i) =
    g_i(\bu)]})^2 \ket{\psi}}.
\end{multline*}
The term inside the first square root is at most 
\begin{equation}\label{eq:add-in-the-bot}
\E_{\bu}  \E_{\bx_1, \dots, \bx_i} \sum_{g_1, \dots, g_i}   \bra{\psi}
  ((\widehat{G}^{\bx_{<i}}_{g_{<i}}  \cdot \widehat{G}^{\bx_i}_{g_i} - \widehat{G}^{\bx_i}_{g_i}  \cdot \widehat{G}^{\bx_{<i}}_{g_{<i}})  \cdot  (\widehat{G}^{\bx_i}_{g_i}\cdot (\widehat{G}^{\bx_{<i}}_{g_{<i}})^\dagger - (\widehat{G}^{\bx_{<i}}_{g_{<i}})^\dagger \cdot \widehat{G}^{\bx_i}_{g_i})) \otimes I \ket{\psi},
\end{equation}
which is at most~$\nu_4$ by \Cref{lem:commute-g-half-sandwich}.
The term inside the second square root is at most~$1$ because~$B$ and~$\widehat{G}$ are measurements.
Next, we claim that
\begin{equation}\label{eq:even-bigger-CS}
\eqref{eq:gonna-need-a-bigger-cauchy-schwarz}
\approx_{\sqrt{\nu_4}}
\E_{\bu}  \E_{\bx_1, \dots, \bx_i} \sum_{g_1, \dots, g_i: g_i \neq \bot} \bra{\psi}\widehat{G}^{\bx_i}_{g_i} \cdot
  \widehat{G}^{\bx_{<i}}_{g_{<i}}  \cdot (\widehat{G}^{\bx_{<i}}_{g_{<i}} )^\dagger \cdot  \widehat{G}^{\bx_i}_{g_i}\ot (I - B^{\bu}_{[f(\bx_i) =
    g_i(\bu)]}) \ket{\psi}.
\end{equation}
To show this, we bound the magnitude of the difference.
\begin{multline*}
\E_{\bu}  \E_{\bx_1, \dots, \bx_i} \sum_{g_1, \dots, g_i: g_i \neq \bot}  \bra{\psi}
   (\widehat{G}^{\bx_i}_{g_i} \cdot \widehat{G}^{\bx_{<i}}_{g_{<i}} \ot (I - B^{\bu}_{[f(\bx_i) =
    g_i(\bu)]})) \cdot (((\widehat{G}^{\bx_{<i}}_{g_{<i}})^\dagger  \cdot \widehat{G}^{\bx_i}_{g_i} - \widehat{G}^{\bx_i}_{g_i}  \cdot (\widehat{G}^{\bx_{<i}}_{g_{<i}})^\dagger) \otimes I) \ket{\psi}\\
    \leq 
    \sqrt{\E_{\bu}  \E_{\bx_1, \dots, \bx_i} \sum_{g_1, \dots, g_i: g_i \neq \bot}  \bra{\psi}
  (\widehat{G}^{\bx_i}_{g_i} \cdot\widehat{G}^{\bx_{<i}}_{g_{<i}}  \cdot  (\widehat{G}^{\bx_{<i}}_{g_{<i}} )^\dagger \cdot \widehat{G}^{\bx_i}_{g_i}) \ot (I - B^{\bu}_{[f(\bx_i) =
    g_i(\bu)]})^2 \ket{\psi}}\\
    \cdot
    \sqrt{\E_{\bu}  \E_{\bx_1, \dots, \bx_i} \sum_{g_1, \dots, g_i: g_i \neq \bot}   \bra{\psi}
  ((\widehat{G}^{\bx_{<i}}_{g_{<i}}  \cdot \widehat{G}^{\bx_i}_{g_i} - \widehat{G}^{\bx_i}_{g_i}  \cdot \widehat{G}^{\bx_{<i}}_{g_{<i}})  \cdot  ((\widehat{G}^{\bx_{<i}}_{g_{<i}})^\dagger  \cdot \widehat{G}^{\bx_i}_{g_i} - \widehat{G}^{\bx_i}_{g_i}  \cdot (\widehat{G}^{\bx_{<i}}_{g_{<i}}))^\dagger) \otimes I \ket{\psi}}.
\end{multline*}
The term inside the first square root is at most~$1$ because~$B$ and $\widehat{G}$ are measurements.
The term inside the second square root is at most \Cref{eq:add-in-the-bot},
which is at most~$\nu_4$ by \Cref{lem:commute-g-half-sandwich}.
But
\begin{align*}
\sum_{g_1, \ldots, g_{i-1}} \widehat{G}^{x_{<i}}_{g_{<i}} \cdot (\widehat{G}^{x_{<i}}_{g_{<i}})^\dagger
&= \sum_{g_1, \ldots, g_{i-1}} \widehat{G}^{x_1}_{g_1} \cdots \widehat{G}^{x_{i-1}}_{g_{i-1}} \cdot \widehat{G}^{x_{i-1}}_{g_{i-1}}  \cdots \widehat{G}^{x_{1}}_{g_{1}}\\
&= \sum_{g_1, \ldots, g_{i-2}} \widehat{G}^{x_1}_{g_1} \cdots \Big(\sum_{g_{i-1}} \widehat{G}^{x_{i-1}}_{g_{i-1}} \Big)  \cdots \widehat{G}^{x_{1}}_{g_{1}}\\
&= \sum_{g_1, \ldots, g_{i-2}} \widehat{G}^{x_1}_{g_1} \cdots I  \cdots \widehat{G}^{x_{1}}_{g_{1}}\\
& \cdots\\
& = I.
\end{align*}
Thus,
\begin{equation*}
\eqref{eq:even-bigger-CS}
= 
\E_{\bu}  \E_{\bx_1, \dots, \bx_i} \sum_{g_i: g_i \neq \bot} \bra{\psi}\widehat{G}^{\bx_i}_{g_i} \ot (I - B^{\bu}_{[f(\bx_i) =
    g_i(\bu)]}) \ket{\psi} \leq \nu_1,
\end{equation*}
by \Cref{eq:ld-gbcon}.
In total, using $\sqrt{426} \leq 21$, this gives an error of
\begin{align*}
\nu_1 + 2 \sqrt{\nu_4}
& = \zeta + \sqrt{8m\eps + 4\delta} + 2\cdot \sqrt{426 k^2 m \cdot \left(\gamma^{1/16} + \zeta^{1/16} +  (d/q)^{1/16}\right)}\\
& \leq \zeta^{1/32} + 3m \eps^{1/32} + 2 \delta^{1/32} + 42 km \cdot \left(\gamma^{1/32} + \zeta^{1/32} + (d/q)^{1/32}\right)\\
& \leq 43 km \cdot \left(\eps^{1/32} + \delta^{1/32} +\gamma^{1/32} + \zeta^{1/32} +  (d/q)^{1/32}\right).
\end{align*}
This completes the proof.
\end{proof}

\subsection{Consistency of~$H$ with~$A$}
\label{sec:consistency-of-h-with-a}

\begin{lemma}[Consistency of~$H$ with~$B$]\label{lem:h-b-consistency}
\begin{equation*}
H_{[h|_u =f]} \otimes I \simeq_{\nu_6} I \otimes B^u_f,
\end{equation*}
where
\begin{equation*}
\nu_6 = 44 k^2m \cdot \left(\eps^{1/32} + \delta^{1/32} + \gamma^{1/32} + \zeta^{1/32} + (d/q)^{1/32}\right).
\end{equation*}
\end{lemma}
\begin{proof}
Let $\bx_1, \ldots, \bx_k \sim \F_q$ be independent and uniformly random.
Our goal is to bound
\begin{align}
&\E_{\bu} \sum_{f \neq f'} \bra{\psi} H_{[h|_{\bu} = f']} \otimes B^{\bu}_f \ket{\psi}\nonumber\\
=~&\E_{\bu} \sum_{h}\sum_{f \neq h|_{\bu}} \bra{\psi} H_{h} \otimes B^{\bu}_f \ket{\psi}\nonumber\\
=~&\E_{\bu}\E_{\bx_1, \ldots, \bx_k} \sum_{h}\sum_{f \neq h|_{\bu}} \bra{\psi} H^{\bx_1, \ldots, \bx_k}_{h} \otimes B^{\bu}_f \ket{\psi}\nonumber\\
=~&\E_{\bu}\E_{\bx_1, \ldots, \bx_k} \sum_{h} \sum_{w: |w| \geq d+1} \sum_{f \neq h|_{\bu}} \bra{\psi} \widehat{H}^{\bx_1, \ldots, \bx_k}_{h_w} \otimes B^{\bu}_f \ket{\psi}\nonumber\\
=~&\E_{\bu}\E_{\bx_1, \ldots, \bx_k} \sum_{h} \sum_{w: |w| \geq d+1} \sum_{(g_1, \ldots, g_k) = h_w} \sum_{f \neq h|_{\bu}} \bra{\psi} \widehat{H}^{\bx_1, \ldots, \bx_k}_{g_1, \ldots, g_k} \otimes B^{\bu}_f \ket{\psi}.\label{eq:keep-on-expandin}
\end{align}
We note that the sum over $(g_1, \ldots, g_k) = h_w$ in the final step is trivial
because there is only ever one tuple $(g_1, \ldots, g_k)$ which is equal to $h_w$.
Because $|w| \geq d+1$,
there exist at least $(d+1)$ coordinates~$i$ such that $g_i \neq \bot$
and hence $g_i = h|_{\bx_i}$.
Since~$f$ is degree-$d$, if it is not equal to $h|_{\bu}$,
then there must exist an $i$ such that $g_i \neq \bot$ and $g_i(\bu) \neq f(\bx_i)$.
Thus,
\begin{align}
\eqref{eq:keep-on-expandin}
~=~& \E_{\bu}\E_{\bx_1, \ldots, \bx_k} \sum_{h} \sum_{w: |w| \geq d+1} \sum_{(g_1, \ldots, g_k) = h_w} \sum_{\substack{f: \exists i: g_i \neq \bot, \\g_i(\bu) \neq f(\bx_i)} }\bra{\psi} \widehat{H}^{\bx_1, \ldots, \bx_k}_{g_1, \ldots, g_k} \otimes B^{\bu}_f \ket{\psi}\nonumber\\
\leq~&\E_{\bu}\E_{\bx_1, \ldots, \bx_k} \sum_{g_1, \ldots, g_k} \sum_{\substack{f: \exists i: g_i \neq \bot, \\g_i(\bu) \neq f(\bx_i)}} \bra{\psi} \widehat{H}^{\bx_1, \ldots, \bx_k}_{g_1, \ldots, g_k} \otimes B^{\bu}_f \ket{\psi}.\label{eq:keep-on-contractin}
\end{align}
Let $(\by_1, \ldots, \by_k) \sim \distinct{k}$. Then by \Cref{prop:ld-dnoteq}, \Cref{eq:keep-on-contractin}
is $(k^2/q)$-close to
\begin{align*}
&\E_{\bu}\E_{\by_1, \ldots, \by_k} \sum_{g_1, \ldots, g_k} \sum_{\substack{f: \exists i: g_i \neq \bot, \\g_i(\bu) \neq f(\by_i)}} \bra{\psi} \widehat{H}^{\by_1, \ldots, \by_k}_{g_1, \ldots, g_k} \otimes B^{\bu}_f \ket{\psi}\\
\leq~&
\sum_i \E_{\bu}\E_{\by_1, \ldots, \by_k} \sum_{g_1, \ldots, g_k} \sum_{\substack{f: g_i \neq \bot, \\g_i(\bu) \neq f(\by_i)}} \bra{\psi} \widehat{H}^{\by_1, \ldots, \by_k}_{g_1, \ldots, g_k} \otimes B^{\bu}_f \ket{\psi} \tag{by the union bound}\\
=~&
\sum_i \E_{\bu}\E_{\by_1, \ldots, \by_k} \sum_{g_1, \ldots, g_k:g_i \neq \bot} \sum_{a \neq g_i(\bu)} \bra{\psi} \widehat{H}^{\by_1, \ldots, \by_k}_{g_1, \ldots, g_k} \otimes B^{\bu}_{[f(\by_i)=a]} \ket{\psi}\\
\leq~& \sum_i \nu_5 \tag{by \Cref{lem:ld-sandwich-line-one-point}}\\
=~& k \cdot \nu_5.
\end{align*}
In total, using $1/q \leq (d/q)^{1/32}$, this gives an error of
\begin{align*}
\frac{k^2}{q} + k \cdot \nu_5
& = \frac{k^2}{q} + k \cdot 43 km \cdot \left(\eps^{1/32} + \delta^{1/32} + \gamma^{1/32} + \zeta^{1/32} + (d/q)^{1/32}\right)\\
& = 44 k^2m \cdot \left(\eps^{1/32} + \delta^{1/32} + \gamma^{1/32} + \zeta^{1/32} + (d/q)^{1/32}\right).
\end{align*}
This completes the proof.
\end{proof}

\begin{corollary}[Consistency of~$H$ with~$A$; Proof of \Cref{item:ld-pasting-N-consistency-sub-measurement} in \Cref{lem:ld-pasting-sub-measurement}]
\begin{equation*}
H_{[h(u, x)=a]} \otimes I \simeq_{\nu} I \otimes A^{u,x}_a.
\end{equation*}
\end{corollary}
\begin{proof}
\Cref{lem:h-b-consistency} implies that
\begin{equation}\label{eq:h-b-consistency-at-a-point}
H_{[h(u, x)=a]} \otimes I \simeq_{\nu_6} I \otimes B^{u}_{[f(x)=a]}.
\end{equation}
\Cref{prop:triangle-sub} applied to \Cref{eq:h-b-consistency-at-a-point} and \Cref{eq:ld-abcon} implies that
\begin{equation*}
H_{[h(u, x)=a]} \otimes I \simeq_{\nu_6 + \sqrt{8m \eps + 4\delta}} I \otimes A^{u,x}_a.
\end{equation*}
We can bound this error by
\begin{align*}
\sqrt{8m \eps + 4\delta} + \nu_6
& = \sqrt{8m \eps + 4\delta}
	+ 44 k^2m \cdot \left(\eps^{1/32} + \delta^{1/32} +  \gamma^{1/32} +\zeta^{1/32} + (d/q)^{1/32}\right)\\
& \leq 3m \eps^{1/32} + 2\delta^{1/32}
	+ 44 k^2m \cdot \left(\eps^{1/32} + \delta^{1/32} +  \gamma^{1/32} + \zeta^{1/32} +(d/q)^{1/32}\right)\\
& \leq  47 k^2m \cdot \left(\eps^{1/32} + \delta^{1/32} +  \gamma^{1/32} + \zeta^{1/32} +(d/q)^{1/32}\right).
\end{align*}
This is clearly at most~$\nu$, which completes the proof.
\end{proof}

\subsection{Completeness of~$H$}
\label{sec:completeness-of-h-low-degree}

\begin{definition}
Let $\tau \in \{0, 1\}^k$ be a type.
We define the following two subsets of $\calP^+(m,q,d)^k$.
\begin{itemize}
\item
We define $\mathsf{Outcomes}_{\tau}$ to be the set of tuples $(g_1, \ldots, g_k)$ such that $g_i  \in \polyfunc{m}{q}{d}$ for each $i \in \tau$ and $g_i = \bot$ for each $i \notin \tau$.
This is the set of possible outcomes of the $\widehat{H}$ measurement of type~$\tau$.
\item
Let $x_1, \ldots, x_k \in \F_q$.
We define $\mathsf{Global}_{\tau}(x)$ to be the subset of $\mathsf{Outcomes}_{\tau}$ containing only those tuples which are consistent with a global polynomial. In other words, $(g_1, \ldots, g_k) \in \mathsf{Global}_{\tau}(x)$ if there exists an $h \in \polyfunc{m+1}{q}{d}$ such that $g_i = h|_{x_i}$ for each $i \in \tau$.
Next, we define
\begin{equation*}
\overline{\mathsf{Global}_{\tau}(x)} = \mathsf{Outcomes}_{\tau} \setminus \mathsf{Global}_{\tau}(x).
\end{equation*}
This contains those tuples of type~$\tau$ with no consistent global polynomial.
\end{itemize}
\end{definition}

\begin{lemma}
\label{lem:over-all-outcomes}
Let $\bx_1, \ldots, \bx_k \sim \F_q$ be sampled uniformly at random.
Then
\begin{equation*}
\bra{\psi} H \otimes I \ket{\psi}
\approx_{\nu_7}  \E_{\bx_1, \ldots, \bx_k} \sum_{\tau:|\tau| \geq d+1} \sum_{(g_1, \ldots, g_k) \in \mathsf{Outcomes}_{\tau}} \bra{\psi} \widehat{H}^{\bx_1, \ldots, \bx_k}_{g_1, \ldots, g_k} \otimes I \ket{\psi},
\end{equation*}
where
\begin{equation*}
\nu_7 = 46 k^2m \cdot \left(\eps^{1/32} + \delta^{1/32} +\gamma^{1/32} +  \zeta^{1/32} + (d/q)^{1/32}\right).
\end{equation*}
\end{lemma}
\begin{proof}
Let $(\by_1, \ldots, \by_k) \sim \mathsf{Distinct}_k$.
By definition,
\begin{align}
\bra{\psi} H \otimes I \ket{\psi}
&= \sum_h \bra{\psi} H_h \otimes I \ket{\psi}\nonumber\\
&= \E_{\by_1, \ldots, \by_k} \sum_h \bra{\psi} H^{\by_1, \ldots, \by_k}_h \otimes I \ket{\psi}\nonumber\\
&= \E_{\by_1, \ldots, \by_k} \sum_h \sum_{\tau:|\tau| \geq d+1} \bra{\psi} \widehat{H}^{\by_1, \ldots, \by_k}_{h_{\tau}} \otimes I \ket{\psi}\nonumber\\
&= \E_{\by_1, \ldots, \by_k} \sum_h \sum_{\tau:|\tau| \geq d+1} \sum_{(g_1, \ldots, g_k) = h_\tau} \bra{\psi} \widehat{H}^{\by_1, \ldots, \by_k}_{g_1, \ldots, g_k} \otimes I \ket{\psi}\nonumber\\
&= \E_{\by_1, \ldots, \by_k}  \sum_{\tau:|\tau| \geq d+1} \sum_{(g_1, \ldots, g_k) \in \mathsf{Global}_{\tau}(\by)} \bra{\psi} \widehat{H}^{\by_1, \ldots, \by_k}_{g_1, \ldots, g_k} \otimes I \ket{\psi}. \label{eq:sum-restricted-to-global-polynomial}
\end{align}
The sum in \Cref{eq:sum-restricted-to-global-polynomial}
is over $g_1, \ldots, g_k$ which are consistent with a global polynomial.
We will now show that the value of this sum
remains largely unchanged if we drop this condition.
In particular, we claim that
\begin{equation}\label{eq:remove-the-restriction}
\eqref{eq:sum-restricted-to-global-polynomial}
\approx_{\frac{k^2}{q} + k \cdot\nu_5 + \frac{md}{q}}
\E_{\by_1, \ldots, \by_k}  \sum_{\tau:|\tau| \geq d+1} \sum_{(g_1, \ldots, g_k) \in \mathsf{Outcomes}_{\tau}} \bra{\psi} \widehat{H}^{\by_1, \ldots, \by_k}_{g_1, \ldots, g_k} \otimes I \ket{\psi}.
\end{equation}
To show this, we note that $\eqref{eq:remove-the-restriction} \geq \eqref{eq:sum-restricted-to-global-polynomial}$. Hence, it suffices to upper bound their difference.
\begin{align}
\eqref{eq:remove-the-restriction} - \eqref{eq:sum-restricted-to-global-polynomial}
&= \E_{\by_1, \ldots, \by_k} \sum_{\tau:|\tau| \geq d+1} \sum_{(g_1, \ldots, g_k) \in \overline{\mathsf{Global}_{\tau}(\by)}} \bra{\psi} \widehat{H}^{\by_1, \ldots, \by_k}_{g_1, \ldots, g_k} \otimes I \ket{\psi}\nonumber\\
&=  \E_{\bu} \E_{\by_1, \ldots, \by_k}  \sum_{\tau:|\tau| \geq d+1} \sum_f \sum_{(g_1, \ldots, g_k) \in \overline{\mathsf{Global}_{\tau}(\by)}} \bra{\psi} \widehat{H}^{\by_1, \ldots, \by_k}_{g_1, \ldots, g_k} \otimes B^{\bu}_f \ket{\psi},\label{eq:B-appears-out-of-thin-air}
\end{align}
because~$B$ is a measurement. Next, we claim that
\begin{equation}\label{eq:add-in-indicator-for-f-and-g}
\eqref{eq:B-appears-out-of-thin-air}
\approx_{\frac{k^2}{q} + k \cdot\nu_5} 
\E_{\bu} \E_{\by_1, \ldots, \by_k}   \sum_{\tau:|\tau| \geq d+1} \sum_f \sum_{(g_1, \ldots, g_k) \in \overline{\mathsf{Global}_{\tau}(\by)}} \bra{\psi} \widehat{H}^{\by_1, \ldots, \by_k}_{g_1, \ldots, g_k} \otimes B^{\bu}_f \ket{\psi} \cdot \bone[\forall i\in \tau, f(\by_i) = g_i(\bu)].
\end{equation}
To show this, we note that $\eqref{eq:B-appears-out-of-thin-air} \geq \eqref{eq:add-in-indicator-for-f-and-g}$.
Thus, it suffices to upper bound their difference $\eqref{eq:B-appears-out-of-thin-air} - \eqref{eq:add-in-indicator-for-f-and-g}$. This is given by
\begin{equation}\label{eq:about-to-swap-x-for-y}
\E_{\bu} \E_{\by_1, \ldots, \by_k}  \sum_{\tau: |\tau|\geq d+1} \sum_f  \sum_{(g_1, \ldots, g_k) \in \overline{\mathsf{Global}_{\tau}(\by)}} \bra{\psi} \widehat{H}^{\by_1, \ldots, \by_k}_{g_1, \ldots, g_k} \otimes B^{\bu}_f \ket{\psi} \cdot \bone[\exists i \in \tau, f(\by_i) \neq g_i(\bu)].
\end{equation}
Recall that $\bx_1, \ldots, \bx_k \sim \F_q$ are sampled independently and uniformly at random.
Then \Cref{prop:ld-dnoteq} implies that \Cref{eq:about-to-swap-x-for-y} is $(k^2/q)$-close to
\begin{align*}
 &\E_{\bu} \E_{\bx_1, \ldots, \bx_k}  \sum_{\tau: |\tau|\geq d+1} \sum_f  \sum_{(g_1, \ldots, g_k) \in \overline{\mathsf{Global}_{\tau}(\bx)}} \bra{\psi} \widehat{H}^{\bx_1, \ldots, \bx_k}_{g_1, \ldots, g_k} \otimes B^{\bu}_f \ket{\psi} \cdot \bone[\exists i \in \tau, f(\bx_i) \neq g_i(\bu)]\\
 \leq~& \E_{\bu} \E_{\bx_1, \ldots, \bx_k}   \sum_{\tau:|\tau| \geq d+1} \sum_f \sum_{(g_1, \ldots, g_k) \in \mathsf{Outcomes}_{\tau}} \bra{\psi} \widehat{H}^{\bx_1, \ldots, \bx_k}_{g_1, \ldots, g_k} \otimes B^{\bu}_f \ket{\psi} \cdot \bone[\exists i \in \tau, f(\bx_i) \neq g_i(\bu)]\\
 \leq~& \E_{\bu} \E_{\bx_1, \ldots, \bx_k}   \sum_{\tau} \sum_f \sum_{(g_1, \ldots, g_k) \in \mathsf{Outcomes}_{\tau}} \bra{\psi} \widehat{H}^{\bx_1, \ldots, \bx_k}_{g_1, \ldots, g_k} \otimes B^{\bu}_f \ket{\psi} \cdot \bone[\exists i \in \tau, f(\bx_i) \neq g_i(\bu)]\\
 \leq~&  \E_{\bu} \E_{\bx_1, \ldots, \bx_k}  \sum_{\tau} \sum_f \sum_{(g_1, \ldots, g_k) \in \mathsf{Outcomes}_{\tau}} \bra{\psi} \widehat{H}^{\bx_1, \ldots, \bx_k}_{g_1, \ldots, g_k} \otimes B^{\bu}_f \ket{\psi} \cdot \Big(\sum_{i\in \tau} \bone[f(\bx_i) \neq g_i(\bu)]\Big)\\
 =~& \sum_i  \E_{\bu} \E_{\bx_1, \ldots, \bx_k}  \sum_{g_1, \ldots, g_k : g_i \neq \bot} \sum_{f: f(\bx_i) \neq g_i(\bu)} \bra{\psi} \widehat{H}^{\bx_1, \ldots, \bx_k}_{g_1, \ldots, g_k} \otimes B^{\bu}_f \ket{\psi}\\
 =~& \sum_i  \E_{\bu} \E_{\bx_1, \ldots, \bx_k}  \sum_{g_1, \ldots, g_k : g_i \neq \bot} \sum_{a \neq g_i(\bu)} \bra{\psi} \widehat{H}^{\bx_1, \ldots, \bx_k}_{g_1, \ldots, g_k} \otimes B^{\bu}_{[f(\bx_i)=a]} \ket{\psi}\\
 \leq~& \sum_i (\nu_5) \tag{by \Cref{lem:ld-sandwich-line-one-point}}\\
 =~& k \cdot \nu_5.
\end{align*}

Returning to \Cref{eq:add-in-indicator-for-f-and-g},
we introduce the  notation $\mathsf{Consistent}_{\tau}(g, \by, \bu)$
to indicate whether there is a consistent degree-$d$ polynomial
interpolating the $g_i$'s along the line in direction~$u$.
In other words,
\begin{equation*}
\mathsf{Consistent}_{\tau}(g, \by, \bu)
= \left\{\begin{array}{rl}
	1 & \text{if $\exists f$ such that $f(\by_i) = g_i(\bu)$ for all $i \in \tau$,}\\
	0 & \text{otherwise.}
	\end{array}\right.
\end{equation*}
Clearly, for any $f$,
\begin{equation*}
\bone[\forall i\in \tau, f(\by_i) = g_i(\bu)] \leq \bone[\mathsf{Consistent}_{\tau}(g, \by, \bu)]
\end{equation*}
because the left-hand side indicates whether~$f$ is the consistent degree-$d$ polynomial
interpolating along direction~$u$.
As a result,
\begin{align}
\eqref{eq:add-in-indicator-for-f-and-g}
\leq~&
\E_{\bu} \E_{\by_1, \ldots, \by_k} \sum_{\tau:|\tau|\geq d+1} \sum_f  \sum_{(g_1, \ldots, g_k) \in \overline{\mathsf{Global}_{\tau}(\by)}} \bra{\psi} \widehat{H}^{\by_1, \ldots, \by_k}_{g_1, \ldots, g_k} \otimes B^{\bu}_f \ket{\psi} \cdot \bone[\mathsf{Consistent}_{\tau}(g,\by,\bu)]\nonumber\\
=~&
\E_{\bu} \E_{\by_1, \ldots, \by_k} \sum_{\tau:|\tau|\geq d+1} \sum_{(g_1, \ldots, g_k) \in \overline{\mathsf{Global}_{\tau}(\by)}} \bra{\psi} \widehat{H}^{\by_1, \ldots, \by_k}_{g_1, \ldots, g_k} \otimes \Big(\sum_f B^{\bu}_f\Big) \ket{\psi} \cdot \bone[\mathsf{Consistent}_{\tau}(g,\by,\bu)]\nonumber\\
=~&
\E_{\bu} \E_{\by_1, \ldots, \by_k}   \sum_{\tau:|\tau|\geq d+1} \sum_{(g_1, \ldots, g_k)  \in \overline{\mathsf{Global}_{\tau}(\by)}} \bra{\psi} \widehat{H}^{\by_1, \ldots, \by_k}_{g_1, \ldots, g_k} \otimes I \ket{\psi} \cdot \bone[\mathsf{Consistent}_{\tau}(g,\by,\bu)]\nonumber\\
=~&
 \E_{\by_1, \ldots, \by_k}   \sum_{\tau:|\tau|\geq d+1} \sum_{(g_1, \ldots, g_k) \in \overline{\mathsf{Global}_{\tau}(\by)}} \bra{\psi} \widehat{H}^{\by_1, \ldots, \by_k}_{g_1, \ldots, g_k} \otimes I \ket{\psi} \cdot \Pr_{\bu}[\mathsf{Consistent}_{\tau}(g,\by,\bu)].\label{eq:consistent-indicator}
\end{align}

Let us now fix $(y_1, \ldots, y_k) \in \mathsf{Distinct}_k$,
a type $\tau$ such that $|\tau| \geq d+1$,
and $(g_1, \ldots, g_k) \in \overline{\mathsf{Global}_{\tau}(y)}$.
Suppose without loss of generality that $\tau_1 = \cdots = \tau_{d+1} = 1$,
so that $g_1, \ldots, g_{d+1} \in \polyfunc{m}{q}{d}$.
Then there is a unique polynomial $h^* \in \polyfunc{m+1}{q}{d}$ which interpolates $g_1, \ldots, g_{d+1}$.
In other words, for all $1 \leq i \leq d+1$,
\begin{equation*}
h^*(u, y_i) = g_i(u).
\end{equation*}
In addition, for any $u$, 
$(h^*)|_u$ is the unique degree-$d$ polynomial
which interpolates $g_1, \ldots, g_{d+1}$ along the line in direction~$u$.
Thus, if there is a consistent degree-$d$ polynomial interpolating all the $g_i$'s along the line in direction~$u$,
then it is $(h^*)|_u$.
In math,
\begin{equation*}
\mathsf{Consistent}_{\tau}(g, y,u) = 1
\qquad \text{if and only if} \qquad
\forall i \in \tau, g_i(u) = h^*(u, y_i).
\end{equation*}
On the other hand, because $(g_1, \ldots, g_k) \in \overline{\mathsf{Global}_{\tau}(y)}$, 
there exists an $i^* \in \tau$ such that $g_{i^*} \neq (h^*)|_{y_{i^*}}$.
Hence,
\begin{equation*}
\Pr_{\bu}[\mathsf{Consistent}_{\tau}(g, y,\bu)]
\leq \Pr_{\bu}[g_{i^*}(\bu) = h^*(\bu, y_{i^*})]
= \Pr_{\bu}[g_{i^*}(\bu) = (h^*)|_{y_{i^*}}(\bu)]
\leq \frac{md}{q},
\end{equation*}
by Schwartz-Zippel.
As a result,
\begin{align*}
\eqref{eq:consistent-indicator}
& \leq 
 \E_{\by_1, \ldots, \by_k}   \sum_{\tau:|\tau|\geq d+1} \sum_{(g_1, \ldots, g_k) \in \overline{\mathsf{Global}_{\tau}(\by)}} \bra{\psi} \widehat{H}^{\by_1, \ldots, \by_k}_{g_1, \ldots, g_k} \otimes I \ket{\psi} \cdot \frac{md}{q}\\
 &\leq
  \E_{\by_1, \ldots, \by_k}   \sum_{g_1, \ldots, g_k} \bra{\psi} \widehat{H}^{\by_1, \ldots, \by_k}_{g_1, \ldots, g_k} \otimes I \ket{\psi} \cdot \frac{md}{q}\\
  & = \frac{md}{q}.
\end{align*}
This establishes \Cref{eq:remove-the-restriction}.
Finally, \Cref{prop:ld-dnoteq} implies that \Cref{eq:remove-the-restriction} is $(k^2/q)$-close to
\begin{equation*}
\E_{\bx_1, \ldots, \bx_k}  \sum_{\tau:|\tau| \geq d+1} \sum_{(g_1, \ldots, g_k) \in \mathsf{Outcomes}_{\tau}} \bra{\psi} \widehat{H}^{\bx_1, \ldots, \bx_k}_{g_1, \ldots, g_k} \otimes I \ket{\psi}.
\end{equation*}
In total, this gives an error of
\begin{align*}
2 \frac{k^2}{q}  + \frac{md}{q}+ k \cdot \nu_5
& = 2 \frac{k^2}{q}  + \frac{md}{q} + k \cdot 43 km \cdot \left(\eps^{1/32} + \delta^{1/32} + \gamma^{1/32} + \zeta^{1/32} +  (d/q)^{1/32}\right)\\
& \leq 3 \frac{k^2 md}{q}  +  43 k^2m \cdot \left(\eps^{1/32} + \delta^{1/32} +\gamma^{1/32} +\zeta^{1/32} +  (d/q)^{1/32}\right)\\
& \leq 3 k^2 m (d/q)^{1/32}  +  43 k^2m \cdot \left(\eps^{1/32} + \delta^{1/32} +\gamma^{1/32} +  \zeta^{1/32} + (d/q)^{1/32}\right)\\
& \leq  46 k^2m \cdot \left(\eps^{1/32} + \delta^{1/32} + \gamma^{1/32} +  \zeta^{1/32} +(d/q)^{1/32}\right).
\end{align*}
This completes the proof of the lemma.
\end{proof}




\begin{lemma}
\label{lem:from-H-to-G}
Let $\bx_1, \ldots, \bx_k \sim \F_q$ be sampled uniformly at random.
Then
\begin{equation*}
\E_{\bx_1, \dots, \bx_k} \sum_{\tau:|\tau| \geq d+1} 
	\sum_{(g_1, \ldots, g_{k}) \in \outc_{\tau}} \bra{\psi} \widehat{H}^{\bx_1, \dots, \bx_k}_{g_1, \dots, g_k} \ot I \ket{\psi}
\approx_{\nu_8}
\sum_{i = d+1}^k \binom{k}{i} \bra{\psi} G^i (I-G)^{k-i} \otimes I \ket{\psi},
\end{equation*}
where
\begin{equation*}
\nu_8 = 46 k m \cdot \left( \gamma^{1/32} + \zeta^{1/32} +(d/q)^{1/32}\right).
\end{equation*}
\end{lemma}
\begin{proof}
We begin by introducing some notation that we will use throughout the proof.
Let $\tau \in \{0, 1\}^k$ be a type.
Then we define
\begin{equation*}
\tau_{< \ell}  = (\tau_1, \ldots, \tau_{\ell-1}) \in \{0, 1\}^{\ell-1},
\quad
\tau_{> \ell}  = (\tau_{\ell+1}, \ldots, \tau_{k}) \in \{0,1\}^{k-\ell},
\end{equation*}
and we define $\tau_{\leq \ell}$ and $\tau_{\geq \ell}$ similarly.
In addition, given $(g_1, \ldots, g_k) \in \calP^+(m, q, d)^{k}$, we define
\begin{equation*}
g_{< \ell}  = (g_1, \ldots, g_{\ell-1}) \in \calP^+(m, q, d)^{\ell-1},
\quad
g_{> \ell}  = (g_{\ell+1}, \ldots, g_{k}) \in \calP^+(m, q, d)^{k-\ell}
\end{equation*}
and we define $g_{\leq \ell}$ and $g_{\geq \ell}$ similarly.
Using this notation, we can write
\begin{equation*}
 \wH^{x_{\geq \ell}}_{g_{\geq \ell}} = \wH^{x_{\ell}, \ldots, x_k}_{g_{\ell}, \ldots, g_k}.
\end{equation*}
Next, we introduce the notation
\begin{equation*}
\wG^{x_{\geq \ell}}_{g_{\geq \ell}} = \wG^{x_{\ell}}_{g_{\ell}} \cdots \wG^{x_k}_{g_k}.
\end{equation*}
This satisfies the recurrence relation
\begin{equation}\label{eq:G-recurrence}
\wG^{x_{\geq \ell}}_{g_{\geq \ell}} = \wG^{x_{\ell}}_{g_{\ell}} \cdot \wG^{x_{> \ell}}_{g_{> \ell}}.
\end{equation}
Furthermore, we can write
\begin{equation}\label{eq:split-H-into-two-Gs}
 \wH^{x_{\geq \ell}}_{g_{\geq \ell}} = (\wG^{x_{\geq \ell}}_{g_{\geq \ell}}) \cdot (\wG^{x_{\geq \ell}}_{g_{\geq \ell}})^\dagger.
\end{equation}
Finally, we will write $\sfO_{\tau}$ as shorthand for $\outc_{\tau}$.



To prove the lemma, we will show that for each $1 \leq \ell \leq k$,
\begin{align}
&\E_{\bx_{\geq \ell}} \sum_{\tau:|\tau| \geq d+1}  \sum_{g_{\geq \ell} \in \sfO_{\tau_{\geq \ell}}}
	\bra{\psi} \widehat{H}^{\bx_{\geq \ell}}_{g_{\geq \ell}} \ot (G^{|\tau_{<\ell}|}\cdot (I-G)^{(\ell-1)-|\tau_{<\ell}|}) \ket{\psi}\nonumber\\
\approx_{2\sqrt{2\zeta} + 2 \sqrt{\nu_4}}~&\E_{\bx_{>\ell}} \sum_{\tau:|\tau| \geq d+1} \sum_{g_{> \ell} \in \sfO_{\tau_{> \ell}}}
	\bra{\psi} \widehat{H}^{\bx_{> \ell}}_{g_{> \ell}} \ot (G^{|\tau_{\leq \ell}|} \cdot (I-G)^{\ell-|\tau_{\leq\ell}|}) \ket{\psi}.
	\label{eq:i-think-this-is-what-i'm-supposed-to-prove}
\end{align}
If we then repeatedly apply \Cref{eq:i-think-this-is-what-i'm-supposed-to-prove} for $\ell = 1, \ldots, k$,
we derive
\begin{align*}
&\E_{\bx_1, \dots, \bx_k} \sum_{\tau:|\tau| \geq d+1} 
	\sum_{(g_1, \ldots, g_{k}) \in \outc_{\tau}} \bra{\psi} \widehat{H}^{\bx_1, \dots, \bx_k}_{g_1, \dots, g_k} \ot I \ket{\psi}\\
=~&\E_{\bx_{\geq 1}} \sum_{\tau:|\tau| \geq d+1} 
	\sum_{g_{\geq 1} \in \sfO_{\tau}} \bra{\psi} \widehat{H}^{\bx_{\geq 1}}_{g_{\geq 1}} \ot I \ket{\psi}\\
\approx_{2\sqrt{2\zeta} + 2 \sqrt{\nu_4}}~&\E_{\bx_{\geq 2}} \sum_{\tau:|\tau| \geq d+1} 
	\sum_{g_{\geq 2} \in \sfO_{\tau_{\geq 2}}} \bra{\psi} \widehat{H}^{\bx_{\geq 2}}_{g_{\geq 2}} \ot (G^{|\tau_{\leq 1}|} \cdot (I-G)^{1-|\tau_{\leq1}|}) \ket{\psi}\\
\approx_{2\sqrt{2\zeta} + 2 \sqrt{\nu_4}}~&\E_{\bx_{\geq 3}} \sum_{\tau:|\tau| \geq d+1} 
	\sum_{g_{\geq 3} \in \sfO_{\tau_{\geq 3}}} \bra{\psi} \widehat{H}^{\bx_{\geq 3}}_{g_{\geq 3}} \ot (G^{|\tau_{\leq 2}|} \cdot (I-G)^{2-|\tau_{\leq2}|}) \ket{\psi}\\
	\cdots&\\
\approx_{2\sqrt{2\zeta} + 2 \sqrt{\nu_4}}~& \sum_{\tau:|\tau| \geq d+1} 
	 \bra{\psi}I  \ot (G^{|\tau|} \cdot (I-G)^{k-|\tau|}) \ket{\psi}\\
=~&\sum_{i = d+1}^k \binom{k}{i} \bra{\psi} I \otimes ( G^i (I-G)^{k-i} ) \ket{\psi}.
\end{align*}
In total, using $\sqrt{426}\leq 21$, this gives an error of
\begin{align*}
k \cdot(2\sqrt{2\zeta} + 2 \sqrt{\nu_4})
& = k \cdot 2 \sqrt{2 \zeta} + 2 \cdot \sqrt{426 k^2 m \cdot \left( \gamma^{1/16} +\zeta^{1/16} + (d/q)^{1/16}\right)}\\
& \leq 4k \cdot \zeta^{1/32} + 42 k m \cdot \left(\gamma^{1/32} + \zeta^{1/32} + (d/q)^{1/32}\right)\\
& \leq 46 k m \cdot \left(\gamma^{1/32} +\zeta^{1/32} +  (d/q)^{1/32}\right).
\end{align*}
This proves the lemma.

We now prove \Cref{eq:i-think-this-is-what-i'm-supposed-to-prove}.
To begin, for each $1 \leq \ell \leq k+1$ and $\tau_{\geq \ell} \in \{0, 1\}^{k - \ell + 1}$,
we define the matrix
\begin{equation}\label{eq:S-def}
S_{\tau_{\geq \ell}} = \sum_{\tau_{< \ell} : |\tau| \geq d+1} G^{|\tau_{<\ell}|}\cdot (I-G)^{(\ell-1)-|\tau_{<\ell}|}.
\end{equation}
Then the statement in \Cref{eq:i-think-this-is-what-i'm-supposed-to-prove} can be rewritten as
\begin{align}
&\E_{\bx_{\geq \ell}} \sum_{\tau_{\geq \ell}}  \sum_{g_{\geq \ell} \in \sfO_{\tau_{\geq \ell}}}
	\bra{\psi} \widehat{H}^{\bx_{\geq \ell}}_{g_{\geq \ell}} \ot S_{\tau_{\geq \ell}}\ket{\psi}\nonumber\\
\approx_{2\sqrt{2\zeta} + 2 \sqrt{\nu_4}}~&\E_{\bx_{> \ell}} \sum_{\tau_{>\ell}} \sum_{g_{>\ell} \in \sfO_{\tau_{> \ell}}}
	\bra{\psi} \widehat{H}^{\bx_{>\ell}}_{g_{>\ell}} \ot S_{\tau_{>\ell}} \ket{\psi}.
	\label{eq:i-think-this-is-what-i'm-supposed-to-prove-2}
\end{align}
To prove this, we will use several facts about $S_{\tau_{\geq \ell}}$.
First, $S$ is Hermitian and positive semidefinite.
This is because each term in \Cref{eq:S-def}
is a product of~$G$ and $(I-G)$.
These matrices commute with each other,
and both are Hermitian and positive semidefinite.
Next, $S$ is bounded:
\begin{align}
S_{\tau_{\geq \ell}}
&= \sum_{\tau_{< \ell} : |\tau| \geq d+1} G^{|\tau_{<\ell}|}\cdot (I-G)^{(\ell-1)-|\tau_{<\ell}|}\nonumber\\
&\leq \sum_{\tau_{< \ell} } G^{|\tau_{<\ell}|}\cdot (I-G)^{(\ell-1)-|\tau_{<\ell}|}\nonumber\\
& = (G + (I-G))^{\ell-1}\nonumber\\
& = I. \label{eq:S-bound}
\end{align}
In addition, for any $\tau_{\ell} \in \{0, 1\}$,
\begin{align}
\Big(\E_{\bx_{\ell}} \sum_{g_{\ell} \in \sfO_{\tau_{\ell}}} \wG^{\bx_{\ell}}_{g_{\ell}}\Big)
& = \left\{ \begin{array}{cl}
		G & \text{if } \tau_{\ell} = 1,\\
		(I-G) & \text{if } \tau_{\ell} = 0,
	\end{array}\right.\nonumber\\
&= G^{\tau_{\ell}} \cdot (I - G)^{1 - \tau_{\ell}}.\label{eq:explicit-formula-for-G-expectation}
\end{align}
Thus, for any $\tau_{> \ell}$,
\begin{align}
\sum_{\tau_{\ell}} 
	S_{\tau_{\geq \ell}} \cdot \Big(\E_{\bx_{\ell}} \sum_{g_{\ell} \in \sfO_{\tau_{\ell}}} \wG^{\bx_{\ell}}_{g_{\ell}}\Big)
&= \sum_{\tau_{\ell}}  S_{\tau_{\geq \ell}} \cdot(G^{\tau_{\ell}} \cdot (I - G)^{1 - \tau_{\ell}})\nonumber\\
&= \sum_{\tau_{\ell}}  \sum_{\tau_{< \ell} : |\tau| \geq d+1} G^{|\tau_{<\ell}|}\cdot (I-G)^{(\ell-1)-|\tau_{<\ell}|}
	\cdot (G^{\tau_{\ell}} \cdot (I - G)^{1 - \tau_{\ell}})\nonumber\\
&= \sum_{\tau_{\ell}}  \sum_{\tau_{< \ell} : |\tau| \geq d+1} G^{|\tau_{\leq \ell}|}\cdot (I-G)^{\ell-|\tau_{\leq\ell}|}\nonumber\\
&= \sum_{\tau_{\leq \ell} : |\tau| \geq d+1} G^{|\tau_{\leq \ell}|}\cdot (I-G)^{\ell-|\tau_{\leq\ell}|}\nonumber\\
&= S_{\tau_{> \ell}}.\label{eq:S-recurrence}
\end{align}
Finally, for any $\tau_{\geq \ell}$,
\begin{align}
&S_{\tau_{\geq \ell}}
	\cdot \Big(\E_{\bx_{\ell}} \sum_{g_{\ell} \in \sfO_{\tau_{\ell}}} \wG^{\bx_{\ell}}_{g_{\ell}}\Big)
	\cdot S_{\tau_{\geq \ell}}\nonumber\\
 =~& S_{\tau_{\geq \ell}}
	\cdot (G^{\tau_{\ell}} \cdot (I - G)^{1 - \tau_{\ell}})
	\cdot S_{\tau_{\geq \ell}} \tag{by \Cref{eq:explicit-formula-for-G-expectation}}\\
 =~& \sqrt{G^{\tau_{\ell}} \cdot (I - G)^{1 - \tau_{\ell}}} \cdot (S_{\tau_{\geq \ell}})^2
	\cdot \sqrt{G^{\tau_{\ell}} \cdot (I - G)^{1 - \tau_{\ell}}}\tag{because $S_{\tau_{\geq \ell}}$ commutes with $G$ and $(I-G)$}\\
\leq~& \sqrt{G^{\tau_{\ell}} \cdot (I - G)^{1 - \tau_{\ell}}} \cdot I
	\cdot \sqrt{G^{\tau_{\ell}} \cdot (I - G)^{1 - \tau_{\ell}}} \tag{by \Cref{eq:S-bound}}\\
=~& G^{\tau_{\ell}} \cdot (I - G)^{1 - \tau_{\ell}}\nonumber\\
=~& \Big(\E_{\bx_{\ell}} \sum_{g_{\ell} \in \sfO_{\tau_{\ell}}} \wG^{\bx_{\ell}}_{g_{\ell}}\Big), \label{eq:S-sandwich}
\end{align}
where the last step uses \Cref{eq:explicit-formula-for-G-expectation} again.
This concludes the set of facts we will need about $S_{\tau_{\geq \ell}}$.

Now we prove \Cref{eq:i-think-this-is-what-i'm-supposed-to-prove-2}.
To start, we write $\wH$ as a sandwich of $\wG$ operators, and move the
rightmost $\wG^{\bx_{\ell}}_{g_{\ell}}$ to the second tensor factor.
\begin{align}
&\E_{\bx_{\geq \ell}} \sum_{\tau_{\geq \ell}}  \sum_{(g_{\ell}, \ldots,g_k) \in \sfO_{\tau_{\geq \ell}}}
	\bra{\psi} \widehat{H}^{\bx_{\ell}, \dots, \bx_{k}}_{g_{\ell}, \dots, g_{k}} \ot S_{\tau_{\geq \ell}}\ket{\psi}\nonumber\\
=~&\E_{\bx_{\geq \ell}} \sum_{\tau_{\geq \ell}}  \sum_{(g_{\ell}, \ldots,g_k) \in \sfO_{\tau_{\geq \ell}}}
	\bra{\psi} (\wG^{\bx_{\geq \ell}}_{g_{\geq \ell}} \cdot (\wG^{\bx_{\geq \ell}}_{g_{\geq \ell}})^\dagger) \ot S_{\tau_{\geq \ell}}\ket{\psi}\tag{by \Cref{eq:split-H-into-two-Gs}}\\
=~&\E_{\bx_{\geq \ell}} \sum_{\tau_{\geq \ell}}  \sum_{(g_{\ell}, \ldots,g_k) \in \sfO_{\tau_{\geq \ell}}}
	\bra{\psi} (\wG^{\bx_{\geq \ell}}_{g_{\geq \ell}} \cdot (\wG^{\bx_{> \ell}}_{g_{> \ell}})^\dagger \cdot \wG^{\bx_{\ell}}_{g_{\ell}}) \ot S_{\tau_{\geq \ell}}\ket{\psi}\tag{by \Cref{eq:G-recurrence}}\\
\approx_{\sqrt{2\zeta}}~&\E_{\bx_{\geq \ell}} \sum_{\tau_{\geq \ell}}  \sum_{(g_{\ell}, \ldots,g_k) \in \sfO_{\tau_{\geq \ell}}}
	\bra{\psi} (\wG^{\bx_{\geq \ell}}_{g_{\geq \ell}} \cdot (\wG^{\bx_{> \ell}}_{g_{> \ell}})^\dagger) \ot (S_{\tau_{\geq \ell}} \cdot \wG^{\bx_{\ell}}_{g_{\ell}})\ket{\psi}.\label{eq:move-g-over-there}
\end{align}
To justify the approximation, we bound the error.
\begin{multline}\label{eq:call-again-later-part-tres}
\Big|\E_{\bx_{\geq \ell}} \sum_{\tau_{\geq \ell}}  \sum_{(g_{\ell}, g_{>\ell}) \in \mathsf{O}_{\tau_{\geq \ell}}}
	\bra{\psi} (\wG^{\bx_{\geq \ell}}_{g_{\geq \ell}}
		\ot S_{\tau_{\geq \ell}}) \cdot (((\wG^{\bx_{> \ell}}_{g_{> \ell}})^\dagger \ot I) \cdot (\wG^{\bx_{\ell}}_{g_{\ell}} \ot I - I \ot \wG^{\bx_{\ell}}_{g_{\ell}} ))\ket{\psi}\Big|\\
\leq \sqrt{\E_{\bx_{\geq \ell}} \sum_{\tau_{\geq \ell}}  \sum_{(g_{\ell}, g_{>\ell}) \in \mathsf{O}_{\tau_{\geq \ell}}}
	\bra{\psi} (\wG^{\bx_{\geq \ell}}_{g_{\geq \ell}} \cdot (\wG^{\bx_{\geq \ell}}_{g_{\geq \ell}})^\dagger)
		\ot (S_{\tau_{\geq \ell}})^2 \ket{\psi}}\\
\cdot \sqrt{\E_{\bx_{\geq \ell}} \sum_{\tau_{\geq \ell}}  \sum_{
	(g_{\ell}, g_{>\ell}) \in \mathsf{O}_{\tau_{\geq \ell}}}
		\bra{\psi} ((\wG^{\bx_{\ell}}_{g_{\ell}} \ot I - I \ot \wG^{\bx_{\ell}}_{g_{\ell}} )
	\cdot (\wG^{\bx_{> \ell}}_{g_{> \ell}} \cdot(\wG^{\bx_{> \ell}}_{g_{> \ell}})^\dagger\ot I)
	\cdot (\wG^{\bx_{\ell}}_{g_{\ell}} \ot I - I \ot \wG^{\bx_{\ell}}_{g_{\ell}} ))\ket{\psi}}.
\end{multline}
The expression inside the first square root is equal to 
\begin{equation*}
\E_{\bx_{\geq \ell}} \sum_{\tau_{\geq \ell}}  \sum_{(g_{\ell}, g_{>\ell}) \in \mathsf{O}_{\tau_{\geq \ell}}}
	\bra{\psi} \wH^{\bx_{\geq \ell}}_{g_{\geq \ell}}
		\ot (S_{\tau_{\geq \ell}})^2 \ket{\psi},
\end{equation*}
which is at most~$1$ because $S_{\tau_{\geq \ell}} \leq I$ and $\widehat{H}$ is a sub-measurement.
The expression inside the second square root is equal to
\begin{align*}
&\E_{\bx_{\geq \ell}} \sum_{\tau_{\geq \ell}}  \sum_{
	(g_{\ell}, g_{>\ell}) \in \mathsf{O}_{\tau_{\geq \ell}}}
		\bra{\psi} ((\wG^{\bx_{\ell}}_{g_{\ell}} \ot I - I \ot \wG^{\bx_{\ell}}_{g_{\ell}} )
	\cdot (\wH^{\bx_{> \ell}}_{g_{> \ell}} \ot I)
	\cdot (\wG^{\bx_{\ell}}_{g_{\ell}} \ot I - I \ot \wG^{\bx_{\ell}}_{g_{\ell}} ))\ket{\psi}\\
\leq~&\E_{\bx_{\geq \ell}} \sum_{g_{\ell}, \ldots, g_{k}}
		\bra{\psi} ((\wG^{\bx_{\ell}}_{g_{\ell}} \ot I - I \ot \wG^{\bx_{\ell}}_{g_{\ell}} )
	\cdot (\wH^{\bx_{> \ell}}_{g_{> \ell}} \ot I)
	\cdot (\wG^{\bx_{\ell}}_{g_{\ell}} \ot I - I \ot \wG^{\bx_{\ell}}_{g_{\ell}} ))\ket{\psi}\\
\leq~& \E_{\bx_{\geq \ell}} \sum_{g_{\ell}}
		\bra{\psi} (\wG^{\bx_{\ell}}_{g_{\ell}} \ot I - I \ot \wG^{\bx_{\ell}}_{g_{\ell}} )^2\ket{\psi}
			\tag{because $\widehat{H}$ is a sub-measurement}\\
\leq~& 2\zeta. \tag{by \Cref{cor:G-hat-facts}}
\end{align*}
We will now commute the leftmost $\wG^{x_\ell}_{g_\ell}$ in
\Cref{eq:move-g-over-there} to the right in two stages. In the first stage, we have:
\begin{align}
\eqref{eq:move-g-over-there}
&= \E_{\bx_{\geq \ell}} \sum_{\tau_{\geq \ell}}  \sum_{(g_{\ell}, g_{>\ell}) \in \sfO_{\tau_{\geq \ell}}}
	\bra{\psi} (\wG^{\bx_{\ell}}_{g_{\ell}} \cdot \wG^{\bx_{> \ell}}_{g_{> \ell}} \cdot (\wG^{\bx_{> \ell}}_{g_{> \ell}})^\dagger) \ot (S_{\tau_{\geq \ell}} \cdot \wG^{\bx_{\ell}}_{g_{\ell}})\ket{\psi}\nonumber\\
& \approx_{\sqrt{\nu_4}} \E_{\bx_{\geq \ell}} \sum_{\tau_{\geq \ell}}  \sum_{(g_{\ell}, g_{>\ell}) \in \sfO_{\tau_{\geq \ell}}}
	\bra{\psi} (\wG^{\bx_{> \ell}}_{g_{> \ell}} \cdot \wG^{\bx_{\ell}}_{g_{\ell}} \cdot (\wG^{\bx_{> \ell}}_{g_{> \ell}})^\dagger) \ot (S_{\tau_{\geq \ell}} \cdot \wG^{\bx_{\ell}}_{g_{\ell}})\ket{\psi}.
\label{eq:commute-g-part-one}
\end{align}
To justify the approximation, we bound the magnitude of the difference.
\begin{multline}\label{eq:call-this-later}
\Big|  \E_{\bx_{\geq \ell}} \sum_{\tau_{\geq \ell}}
  \sum_{(g_{\ell}, g_{>\ell}) \in \sfO_{\tau_{\geq \ell}}}
  \bra{\psi}     ([ \wG^{\bx_{>\ell}}_{g_{>\ell}},
    \wG^{\bx_\ell}_{g_\ell} ] \otimes I) \cdot     ((\wG^{\bx_{>\ell}}_{g_{>\ell}})^\dagger \ot
    (S_{\tau_{\geq \ell}} \cdot
  \wG^{\bx_\ell}_{g_\ell})) \ket{\psi} \Big|  \\
  \leq \sqrt{ \E_{\bx_{\geq \ell}} \sum_{\tau_{\geq \ell}}
  \sum_{(g_{\ell}, g_{>\ell}) \in \sfO_{\tau_{\geq \ell}}} \bra{\psi} ([ \wG^{\bx_{>\ell}}_{g_{>\ell}},
    \wG^{\bx_\ell}_{g_\ell} ]) \cdot ([ \wG^{\bx_{>\ell}}_{g_{>\ell}},
    \wG^{\bx_\ell}_{g_\ell} ])^\dagger \ot I \ket{\psi}}  \\
  \quad \cdot \sqrt{\E_{\bx_{\geq \ell}} \sum_{\tau_{\geq \ell}}
   \sum_{(g_{\ell}, g_{>\ell}) \in \sfO_{\tau_{\geq \ell}}} \bra{\psi}
    (\wG^{\bx_{>\ell}}_{g_{>\ell}}\cdot (
    \wG^{\bx_{>\ell}}_{g_{>\ell}})^\dagger) \ot
    (\wG^{\bx_{\ell}}_{g_\ell}  \cdot (S_{\tau_{\geq \ell}})^2
    \cdot \wG^{\bx_{\ell}}_{g_\ell}) \ket{\psi}}.
\end{multline}
The quantity inside the first square root is at most
\begin{equation*}
\E_{\bx_{\geq \ell}} \sum_{g_{\ell}, \ldots, g_k} \bra{\psi} ([ \wG^{\bx_{>\ell}}_{g_{>\ell}},
    \wG^{\bx_\ell}_{g_\ell} ]) \cdot ([ \wG^{\bx_{>\ell}}_{g_{>\ell}},
    \wG^{\bx_\ell}_{g_\ell} ])^\dagger \ot I \ket{\psi},
\end{equation*}
which is at most~$\nu_4$ by~\Cref{lem:commute-g-half-sandwich}.
The quantity inside the second square root is equal to
\begin{align*}
&\E_{\bx_{\geq \ell}} \sum_{\tau_{\geq \ell}}
   \sum_{(g_{\ell}, g_{>\ell}) \in \sfO_{\tau_{\geq \ell}}} \bra{\psi}
    \wH^{\bx_{>\ell}}_{g_{>\ell}} \ot
    (\wG^{\bx_{\ell}}_{g_\ell}  \cdot (S_{\tau_{\geq \ell}})^2
    \cdot \wG^{\bx_{\ell}}_{g_\ell}) \ket{\psi}\\
\leq~
&\E_{\bx_{\geq \ell}} \sum_{\tau_{\geq \ell}}
   \sum_{(g_{\ell}, g_{>\ell}) \in \sfO_{\tau_{\geq \ell}}} \bra{\psi}
    \wH^{\bx_{>\ell}}_{g_{>\ell}} \ot
    (\wG^{\bx_{\ell}}_{g_\ell}  \cdot I
    \cdot \wG^{\bx_{\ell}}_{g_\ell}) \ket{\psi} \tag{because $S_{\tau_{\geq \ell}} \leq I$}\\
=~
&\E_{\bx_{\geq \ell}} \sum_{\tau_{\geq \ell}}
   \sum_{(g_{\ell}, g_{>\ell}) \in \sfO_{\tau_{\geq \ell}}} \bra{\psi}
    \wH^{\bx_{>\ell}}_{g_{>\ell}} \ot
    \wG^{\bx_{\ell}}_{g_\ell} \ket{\psi} \\
\leq~&  1. \tag{because $\wG$ and $\wH$ are sub-measurements}
\end{align*}
We continue commuting the leftmost $\wG^{x_{\ell}}_{g_\ell}$ to the
right. 
\begin{align}
  \eqref{eq:commute-g-part-one} =&
  \E_{\bx_{\geq \ell}} \sum_{\tau_{\geq \ell}}  \sum_{(g_{\ell}, g_{>\ell}) \in \sfO_{\tau_{\geq \ell}}}
	\bra{\psi} (\wG^{\bx_{> \ell}}_{g_{> \ell}} \cdot \wG^{\bx_{\ell}}_{g_{\ell}} \cdot (\wG^{\bx_{> \ell}}_{g_{> \ell}})^\dagger) \ot (S_{\tau_{\geq \ell}} \cdot \wG^{\bx_{\ell}}_{g_{\ell}})\ket{\psi} \nonumber\\
  \approx_{\sqrt{\nu_4}}~& 
   \E_{\bx_{\geq \ell}} \sum_{\tau_{\geq \ell}}  \sum_{(g_{\ell}, g_{>\ell}) \in \sfO_{\tau_{\geq \ell}}}
	\bra{\psi} (\wG^{\bx_{> \ell}}_{g_{> \ell}} \cdot (\wG^{\bx_{> \ell}}_{g_{> \ell}})^\dagger \cdot \wG^{\bx_{\ell}}_{g_{\ell}}) \ot (S_{\tau_{\geq \ell}} \cdot \wG^{\bx_{\ell}}_{g_{\ell}})\ket{\psi}. \label{eq:commute-g-part-two}
\end{align}
To justify the approximation, we will need to be slightly more
clever this time. First, as always, we bound the magnitude of the
difference:
\begin{multline}\label{eq:call-again-later-part-dos}
  \Big|  \E_{\bx_{\geq \ell}} \sum_{\tau_{\geq \ell}}  \sum_{(g_{\ell}, g_{>\ell}) \in \sfO_{\tau_{\geq \ell}}}
  \bra{\psi} (\wG^{\bx_{>\ell}}_{g_{>\ell}}         \ot
    (S_{\tau_{\geq \ell}} \cdot
  \wG^{\bx_\ell}_{g_\ell})) \cdot ([ (\wG^{\bx_{>\ell}}_{g_{>\ell}})^\dagger,
    \wG^{\bx_\ell}_{g_\ell} ] \ot I) \ket{\psi} \Big|\\
  \leq   \sqrt{\E_{\bx_{\geq \ell}} \sum_{\tau_{\geq \ell}}  \sum_{(g_{\ell}, g_{>\ell}) \in \sfO_{\tau_{\geq \ell}}}
    \bra{\psi}
    (\wG^{\bx_{>\ell}}_{g_{>\ell}})\cdot (
    \wG^{\bx_{>\ell}}_{g_{>\ell}})^\dagger \ot (S_{\tau_{\geq \ell}} \cdot
    \wG^{\bx_{\ell}}_{g_\ell}  \cdot S_{\tau_{\geq \ell}})
    \ket{\psi}} \\
  \cdot \sqrt{ \E_{\bx_{\geq \ell}}
  \sum_{\tau_{\geq \ell}}  \sum_{(g_{\ell}, g_{>\ell}) \in \sfO_{\tau_{\geq \ell}}}
   \bra{\psi} ([ (\wG^{\bx_{>\ell}}_{g_{>\ell}})^\dagger,
    \wG^{\bx_\ell}_{g_\ell} ])^\dagger \cdot ([ (\wG^{\bx_{>\ell}}_{g_{>\ell}})^\dagger,
    \wG^{\bx_\ell}_{g_\ell} ]) \ot I \ket{\psi}}.
\end{multline}
The term inside the first square root is equal to
\begin{align*}
&\E_{\bx_{\geq \ell}} \sum_{\tau_{\geq \ell}}  \sum_{(g_{\ell}, g_{>\ell}) \in \sfO_{\tau_{\geq \ell}}}
    \bra{\psi}
    \wH^{\bx_{>\ell}}_{g_{>\ell}}
     \ot (S_{\tau_{\geq \ell}} \cdot
    \wG^{\bx_{\ell}}_{g_\ell}  \cdot S_{\tau_{\geq \ell}})
    \ket{\psi}\\
=~&\E_{\bx_{> \ell}} \sum_{\tau_{\geq \ell}}  \sum_{g_{>\ell} \in \sfO_{\tau_{> \ell}}}
    \bra{\psi}
    \wH^{\bx_{>\ell}}_{g_{>\ell}}
     \ot (S_{\tau_{\geq \ell}} \cdot
    \Big(\E_{\bx_{\ell}} \sum_{g_{\ell} \in \sfO_{\tau_{\ell}}}\wG^{\bx_{\ell}}_{g_\ell}  \Big)\cdot S_{\tau_{\geq \ell}})
    \ket{\psi}\\
\leq~&\E_{\bx_{> \ell}} \sum_{\tau_{\geq \ell}}  \sum_{g_{>\ell} \in \sfO_{\tau_{> \ell}}}
    \bra{\psi}
    \wH^{\bx_{>\ell}}_{g_{>\ell}}
     \ot 
    \Big(\E_{\bx_{\ell}} \sum_{g_{\ell} \in \sfO_{\tau_{\ell}}}\wG^{\bx_{\ell}}_{g_\ell}  \Big)
    \ket{\psi} \tag{by~\Cref{eq:S-sandwich}}\\
=~&\E_{\bx_{\geq \ell}} \sum_{\tau_{\geq \ell}}  \sum_{(g_{\ell}, g_{>\ell}) \in \sfO_{\tau_{\geq \ell}}}
    \bra{\psi}
    \wH^{\bx_{>\ell}}_{g_{>\ell}}
     \ot
    \wG^{\bx_{\ell}}_{g_\ell}  
    \ket{\psi}\\
\leq~& 1. \tag{because $\wG$ and $\wH$ are sub-measurements}
\end{align*}
The term inside the second square root is equal to the term inside the first square root of~\Cref{eq:call-this-later}, which we bounded by~$\nu_4$.
We are now ready for the final step, which is to bring the leftmost
$\wG^{x_\ell}_{g_\ell}$ over to the second tensor factor.
\begin{align}
  \eqref{eq:commute-g-part-two} &= 
  \E_{\bx_{\geq \ell}} \sum_{\tau_{\geq \ell}}  \sum_{(g_{\ell}, g_{>\ell}) \in \sfO_{\tau_{\geq \ell}}}
	\bra{\psi} (\wG^{\bx_{> \ell}}_{g_{> \ell}} \cdot (\wG^{\bx_{> \ell}}_{g_{> \ell}})^\dagger \cdot \wG^{\bx_{\ell}}_{g_{\ell}}) \ot (S_{\tau_{\geq \ell}} \cdot \wG^{\bx_{\ell}}_{g_{\ell}})\ket{\psi}\nonumber\\
  &\approx_{\sqrt{2\zeta}} \E_{\bx_{\geq \ell}} \sum_{\tau_{\geq \ell}}  \sum_{(g_{\ell}, g_{>\ell}) \in \sfO_{\tau_{\geq \ell}}}
	\bra{\psi} (\wG^{\bx_{> \ell}}_{g_{> \ell}} \cdot (\wG^{\bx_{> \ell}}_{g_{> \ell}})^\dagger) \ot (S_{\tau_{\geq \ell}} \cdot \wG^{\bx_{\ell}}_{g_{\ell}}\cdot \wG^{\bx_{\ell}}_{g_{\ell}})\ket{\psi}. \label{eq:h-ot-mgg}
\end{align}
To justify the approximation, we bound the magnitude of the
difference.
\begin{multline*}
  \Big|\E_{\bx_{\geq \ell}} \sum_{\tau_{\geq \ell}}
   \sum_{(g_{\ell}, g_{>\ell}) \in \sfO_{\tau_{\geq \ell}}}
   	\bra{\psi} (\wG^{\bx_{> \ell}}_{g_{> \ell}}
		\ot (S_{\tau_{\geq \ell}}
		\cdot \wG^{\bx_{\ell}}_{g_{\ell}}))
		\cdot
		(((\wG^{\bx_{> \ell}}_{g_{> \ell}})^\dagger \ot I)
		\cdot (\wG^{\bx_{\ell}}_{g_{\ell}} \ot I - I \ot \wG^{\bx_{\ell}}_{g_{\ell}}))\ket{\psi}
  \Big|  \\
  \leq \sqrt{\E_{\bx_{\geq \ell}}\sum_{\tau_{\geq \ell}}
   \sum_{(g_{\ell}, g_{>\ell}) \in \sfO_{\tau_{\geq \ell}}}
   	\bra{\psi} (\wG^{\bx_{> \ell}}_{g_{> \ell}} \cdot (\wG^{\bx_{> \ell}}_{g_{> \ell}})^\dagger)
		\ot (S_{\tau_{\geq \ell}}
		\cdot \wG^{\bx_{\ell}}_{g_{\ell}} \cdot S_{\tau_{\geq \ell}}) \ket{\psi}} \nonumber \\
  \cdot \sqrt{\E_{\bx_{\geq \ell}} \sum_{\tau_{\geq \ell}}
   \sum_{(g_{\ell}, g_{>\ell}) \in \sfO_{\tau_{\geq \ell}}} \bra{\psi}
   ( (\wG^{\bx_\ell}_{g_\ell} \ot I - I \ot \wG^{\bx_\ell}_{g_\ell})\cdot
   (\wG^{\bx_{> \ell}}_{g_{> \ell}} \cdot (\wG^{\bx_{> \ell}}_{g_{> \ell}})^\dagger \ot I)
    \cdot (\wG^{\bx_\ell}_{g_\ell} \ot I - I \ot \wG^{\bx_\ell}_{g_\ell})).
    \ket{\psi}}
\end{multline*}
The expression inside the first is equal to the expression inside the first square root in~\Cref{eq:call-again-later-part-dos},
which we bounded by~$1$.
The expression inside the second square root is equal to the expression inside the second square root in~\Cref{eq:call-again-later-part-tres},
which we bounded by~$2\zeta$.
We end by noting that
\begin{align*}
\eqref{eq:h-ot-mgg}
& =
\E_{\bx_{\geq \ell}} \sum_{\tau_{\geq \ell}}  \sum_{(g_{\ell}, g_{>\ell}) \in \sfO_{\tau_{\geq \ell}}}
	\bra{\psi} (\wG^{\bx_{> \ell}}_{g_{> \ell}} \cdot (\wG^{\bx_{> \ell}}_{g_{> \ell}})^\dagger) \ot (S_{\tau_{\geq \ell}} \cdot \wG^{\bx_{\ell}}_{g_{\ell}})\ket{\psi} \tag{because $\wG$ is projective}\\
& =
\E_{\bx_{\geq \ell}} \sum_{\tau_{\geq \ell}}  \sum_{(g_{\ell}, g_{>\ell}) \in \sfO_{\tau_{\geq \ell}}}
	\bra{\psi} \wH^{\bx_{> \ell}}_{g_{> \ell}} \ot (S_{\tau_{\geq \ell}} \cdot \wG^{\bx_{\ell}}_{g_{\ell}})\ket{\psi}\\
& =\E_{\bx_{> \ell}} \sum_{\tau_{> \ell}}  \sum_{g_{>\ell} \in \sfO_{\tau_{> \ell}}}
	\bra{\psi} \wH^{\bx_{> \ell}}_{g_{> \ell}} \ot \Big(\sum_{\tau_{\ell}}S_{\tau_{\geq \ell}} \cdot
		\Big(\E_{\bx_{\ell}} \sum_{g_\ell \in \sfO_{\tau_{\ell}}}\wG^{\bx_{\ell}}_{g_{\ell}}\Big)\Big)\ket{\psi}\\
& =\E_{\bx_{> \ell}} \sum_{\tau_{> \ell}}  \sum_{g_{>\ell} \in \sfO_{\tau_{> \ell}}}
	\bra{\psi} \wH^{\bx_{> \ell}}_{g_{> \ell}} \ot S_{\tau_{>\ell}}\ket{\psi}. \tag{by \Cref{eq:S-recurrence}}
\end{align*}
This concludes the proof of
\Cref{eq:i-think-this-is-what-i'm-supposed-to-prove-2}
and therefore proves the lemma.
\end{proof}



\begin{lemma}
  \label{lem:chernoff-bernoulli-matrix}
  Let $0 < \theta < 1$ and let $k, d > 0$ be integers such that $k \geq 2d/\theta$.
  Define the matrix-valued function $F$ by
  \[ F(X) = \sum_{r = d+1}^{k} \binom{k}{r} X^r (I - X)^{r - k}. \]
  Then for any Hermitian matrix $X$ such that $0 \leq X \leq I$ and $\bra{\psi} X \ot I
  \ket{\psi} \geq 1 - \kappa$, it holds that
  \[ \bra{\psi} F(X) \ot I \ket{\psi} \geq 1 - \frac{\kappa}{1 - \theta}
    - e^{-\theta^2 k/2}. \]
  \end{lemma}
  \begin{proof}
    Let $\rho$ be the reduced state of $\ket{\psi}$ on one prover's
    subsystem (since the state is assumed to be symmetric, it does not
    matter which prover we take).
    Write the eigendecomposition $X = \sum_{i} \lambda_i
    \ket{v_i}\bra{v_i}$ of $X$, where we allow some of the
    eigenvalues to be $0$ so that the set of eigenvectors
    $\{\ket{v_i}\}$ forms an orthonormal basis of the space. This defines a probability
    distribution $\mu$ over eigenvectors, where eigenvector $i$
    occurs with probability $\mu(i) = \bra{v_i} \rho \ket{v_i}$. The
    given condition $\bra{\psi} X \ot I \ket{\psi} = \Tr(X \rho) \geq
    1 -\kappa$ implies that
    \begin{align*}
      \E_{\bi \sim \mu} \lambda_{\bi}
      = \sum_i \lambda_i \cdot \mu(i)
       &= \sum_i \lambda_i \cdot \bra{v_i} \rho \ket{v_i}\\
       &=\sum_i \lambda_i \cdot \tr(\ket{v_i} \bra{v_i}  \cdot \rho)\\
       &=\tr \Big(\sum_i \lambda_i  \ket{v_i} \bra{v_i} \cdot\rho\Big)
       = \tr(X \rho)
       \geq 1 - \kappa.
     \end{align*}
     Or, equivalently,
     \begin{equation*}
      \E_{\bi \sim \mu} ( 1 -\lambda_{\bi}) \leq \kappa.
    \end{equation*}
      By Markov's inequality, for any $0 < \theta < 1$, we have
     \begin{equation*}
     \Pr_{\bi \sim \mu}[(1-\lambda_{\bi}) \geq (1 - \theta)]
     \leq \frac{\E_{\bi \sim \mu} (1-\lambda_{\bi})}{1-\theta}
     \leq \frac{\kappa}{1-\theta}.
     \end{equation*} 
     In other words,
     \begin{equation}\label{eq:in-other-words}
     \Pr_{\bi \sim \mu}[\lambda_{\bi} > \theta]
     = \Pr_{\bi \sim \mu}[(1 - \lambda_{\bi}) < (1 - \theta)]
     = 1 - \Pr_{\bi \sim \mu}[(1 - \lambda_{\bi}) \geq (1 - \theta)]
     \geq 1 - \frac{\kappa}{1-\theta}.  
     \end{equation}
      Thus, it  holds that with probability at least $1 - \frac{\kappa}{1 - \theta}$ over $\bi \sim
    \mu$, $\lambda_{\bi} > \theta$.

    We now evaluate $\bra{\psi} F(X) \ot I \ket{\psi}$. We will
    essentially do this eigenvalue by eigenvalue.
    To begin, we consider a hypothetical eigenvalue $d/k \leq p \leq 1$.
    Observe that $F(p)$ is precisely the
    probability 
    probability of observing at least $d+1$ successes out of $k$
    i.i.d.\ Bernoulli trials, each of which succeeds with probability
    $p$.
    In other words, it is the probability that $\bY := \bY_1 + \cdots + \bY_k \geq d+1$,
    where $\bY_1, \ldots, \bY_k \sim \mathrm{Bernoulli}(p)$.
    We can bound this probability by the additive Chernoff bound (see the second additive bound in~\cite{Blu11}):
    \begin{align*}
    \Pr[\bY \leq d]
    &= \Pr[\bY \leq p k - (pk - d)] \\
    &=\Pr\Big[\bY \leq pk - \Big(p - \frac{d}{k}\Big)\cdot k\Big]\\
    	&\leq  \exp\Big(-2 \Big(p - \frac{d}{k}\Big)^2 \cdot k\Big).
    \end{align*}
    Thus,
    \begin{equation}\label{eq:by-chernoff}
    F(p) = \Pr[\bY \geq d+1] = 1 - \Pr[\bY \leq d] \geq 1 - \exp\Big(-2 \Big(p - \frac{d}{k}\Big)^2 \cdot k\Big).
    \end{equation}
        Putting the pieces together, we compute $\bra{\psi} f(X) \ot I
    \ket{\psi}$:
    \begin{align}
      \bra{\psi} F(X) \ot I \ket{\psi} &= \Tr(F(X) \rho) \nonumber\\
                                       &= \sum_{i} F(\lambda_i)
                                         \bra{v_i} \rho \ket{v_i} \nonumber\\
                                       &=\E_{\bi \sim \mu}
                                         F(\lambda_{\bi}) \nonumber\\
                                       &\geq \Pr_{\bi \sim
                                         \mu}[\lambda_{\bi} \geq \theta]
                                         \cdot F(\theta) \nonumber\\
                                       &\geq \Big( 1 - \frac{\kappa}{1 - \theta} \Big) \cdot F(\theta)  \tag{by \Cref{eq:in-other-words}}\\
      &\geq \Big( 1 - \frac{\kappa}{1 - \theta} \Big) \cdot \Big(1 - \exp \Big(-2\Big(\theta -
        \frac{d}{k}\Big)^2 \cdot k \Big) \Big) \label{eq:almost-done-with-this-giant-proof},
        \end{align}
        where the last step uses \Cref{eq:by-chernoff}.
    Next, we claim that if $a, b, c \geq 0$ satisfy $a \geq (1-b) \cdot (1-c)$,
    then $a \geq 1 - b - c$. This is because if either $b$ or $c$ is at least~$1$,
    then the conclusion is trivially true, and if both are less than~$1$,
    then
    \begin{equation*}
    (1-b) \cdot (1-c) = 1\cdot (1-c) -  b \cdot (1-c) \geq 1 \cdot (1-c) - b \cdot 1 = 1-c-b.
    \end{equation*}
    Since $\bra{\psi} F(X) \ot I \ket{\psi}$ is manifestly positive,
    we can apply this to \Cref{eq:almost-done-with-this-giant-proof}, yielding
    \begin{equation*}
    \bra{\psi} F(X) \ot I \ket{\psi}
    \geq  1 - \frac{\kappa}{1 - \theta} - \exp \Big(-2\Big(\theta -
        \frac{d}{k}\Big)^2 \cdot k \Big).
    \end{equation*}
    Finally, we note that because $k \geq 2d/\theta$,
    we have $\theta/2 \geq d/k$.
    This implies that
    $\theta-d/k \geq \theta - \theta/2 = \theta/2$,
    and so    
    $(\theta - d/k)^2 \geq (\theta/2)^2 = \theta^2/4$.
    As a result,
    \begin{equation*}
    \exp \Big(2\Big(\theta -
        \frac{d}{k}\Big)^2 \cdot k \Big)
        \geq 
        \exp \Big(\theta^2 k/2 \Big).
    \end{equation*}
    Equivalently,
    \begin{equation*}
    \exp \Big(-2\Big(\theta -
        \frac{d}{k}\Big)^2 \cdot k \Big)
        \leq 
        \exp \Big(-\theta^2 k/2 \Big).
    \end{equation*}
    Thus, we conclude
    \begin{equation*}
     \bra{\psi} F(X) \ot I \ket{\psi}
    \geq  1 - \frac{\kappa}{1 - \theta} - \exp \Big(-\theta^2 k/2 \Big).\qedhere
    \end{equation*}
  \end{proof}
\begin{corollary}[Completeness of~$H$; Proof of \Cref{item:ld-pasting-N-completeness-sub-measurement} of \Cref{lem:ld-pasting-sub-measurement}]
  \label{cor:ld-pasting-N-completeness}
  Let $k \geq 400md$. Then
  \[ \bra{\psi} H \otimes I \ket{\psi}
\geq 1 - \kappa \cdot \left(1 + \frac{1}{100m}\right)
    - \nu - e^{- k/(80000m^2)}. \]
\end{corollary}
\begin{proof}
We begin by approximating the completeness as follows.
\begin{align*}
\bra{\psi} H \otimes I \ket{\psi}
& \approx_{\nu_7} \E_{\bx_1, \ldots, \bx_k} \sum_{\tau:|\tau| \geq d+1} \sum_{(g_1, \ldots, g_k) \in \mathsf{Outcomes}_{\tau}} \bra{\psi} \widehat{H}^{\bx_1, \ldots, \bx_k}_{g_1, \ldots, g_k} \otimes I \ket{\psi} \tag{by \Cref{lem:over-all-outcomes}}\\
& \approx_{\nu_8} \sum_{i = d+1}^k \binom{k}{i} \bra{\psi} G^i (I-G)^{k-i} \otimes I \ket{\psi} \tag{by \Cref{lem:from-H-to-G}}\\
& = \bra{\psi} F(G) \ot I \ket{\psi}.
\end{align*}
We can bound the error incurred here by
\begin{align*}
\nu_7 + \nu_8
&= 46 k^2m \cdot \left(\eps^{1/32} + \delta^{1/32} +  \gamma^{1/32} + \zeta^{1/32} +(d/q)^{1/32}\right)
	+ 46 k m \cdot \left(\gamma^{1/32} + \zeta^{1/32} + (d/q)^{1/32}\right)\\
&\leq 100 k^2m \cdot \left(\eps^{1/32} + \delta^{1/32} +  \gamma^{1/32} + \zeta^{1/32} +(d/q)^{1/32}\right)\\
&= \nu.
\end{align*}
Let $\theta = 1/(200 m)$. Note that
\begin{equation*}
\frac{1}{1-\theta}
= \frac{1}{1- 1/(200m)}
= \frac{200m}{200m-1}
= 1 + \frac{1}{200m-1}
\leq 1 + \frac{1}{100m}.
\end{equation*}
Then $k \geq 2d/\theta$ and, therefore, $k \geq d+1$.
As a result, $k$ and~$\theta$ satisfy the hypothesis of \Cref{lem:chernoff-bernoulli-matrix},
namely that $k \geq \max\{d+1, 2d/\theta\}$.
Thus, \Cref{lem:chernoff-bernoulli-matrix} implies that
\begin{align*}
\bra{\psi} F(G) \ot I \ket{\psi}
&\geq 1 - \frac{\kappa}{1 - \theta}
    - e^{-\theta^2 k/2}\\
&= 1 - \frac{\kappa}{1 - \theta}
    - e^{- k/(80000m^2)}\\
&\geq 1 - \kappa \cdot \left(1 + \frac{1}{100m}\right)
    - e^{- k/(80000m^2)}.
\end{align*}
In total, we have
\begin{equation*}
\bra{\psi} H \otimes I \ket{\psi}
\geq 1 - \kappa \cdot \left(1 + \frac{1}{100m}\right)
    - \nu - e^{- k/(80000m^2)}.
\end{equation*}
This completes the proof.
\end{proof}
%%% Local Variables:
%%% mode: latex
%%% TeX-master: "multilinearity"
%%% End:
